% Options for packages loaded elsewhere
% Options for packages loaded elsewhere
\PassOptionsToPackage{unicode}{hyperref}
\PassOptionsToPackage{hyphens}{url}
%
\documentclass[
  ignorenonframetext,
  aspectratio=169,
  spanish,
]{beamer}
\newif\ifbibliography
\usepackage{pgfpages}
\setbeamertemplate{caption}[numbered]
\setbeamertemplate{caption label separator}{: }
\setbeamercolor{caption name}{fg=normal text.fg}
\beamertemplatenavigationsymbolshorizontal
% remove section numbering
\setbeamertemplate{part page}{
  \centering
  \begin{beamercolorbox}[sep=16pt,center]{part title}
    \usebeamerfont{part title}\insertpart\par
  \end{beamercolorbox}
}
\setbeamertemplate{section page}{
  \centering
  \begin{beamercolorbox}[sep=12pt,center]{section title}
    \usebeamerfont{section title}\insertsection\par
  \end{beamercolorbox}
}
\setbeamertemplate{subsection page}{
  \centering
  \begin{beamercolorbox}[sep=8pt,center]{subsection title}
    \usebeamerfont{subsection title}\insertsubsection\par
  \end{beamercolorbox}
}
% Prevent slide breaks in the middle of a paragraph
\widowpenalties 1 10000
\raggedbottom
\AtBeginPart{
  \frame{\partpage}
}
\AtBeginSection{
  \ifbibliography
  \else
    \frame{\sectionpage}
  \fi
}
\AtBeginSubsection{
  \frame{\subsectionpage}
}
\usepackage{iftex}
\ifPDFTeX
  \usepackage[T1]{fontenc}
  \usepackage[utf8]{inputenc}
  \usepackage{textcomp} % provide euro and other symbols
\else % if luatex or xetex
  \usepackage{unicode-math} % this also loads fontspec
  \defaultfontfeatures{Scale=MatchLowercase}
  \defaultfontfeatures[\rmfamily]{Ligatures=TeX,Scale=1}
\fi
\usepackage{lmodern}

\usetheme[]{default}
\usefonttheme[]{default}
\ifPDFTeX\else
  % xetex/luatex font selection
\fi
% Use upquote if available, for straight quotes in verbatim environments
\IfFileExists{upquote.sty}{\usepackage{upquote}}{}
\IfFileExists{microtype.sty}{% use microtype if available
  \usepackage[]{microtype}
  \UseMicrotypeSet[protrusion]{basicmath} % disable protrusion for tt fonts
}{}
\makeatletter
\@ifundefined{KOMAClassName}{% if non-KOMA class
  \IfFileExists{parskip.sty}{%
    \usepackage{parskip}
  }{% else
    \setlength{\parindent}{0pt}
    \setlength{\parskip}{6pt plus 2pt minus 1pt}}
}{% if KOMA class
  \KOMAoptions{parskip=half}}
\makeatother

\usepackage{color}
\usepackage{fancyvrb}
\newcommand{\VerbBar}{|}
\newcommand{\VERB}{\Verb[commandchars=\\\{\}]}
\DefineVerbatimEnvironment{Highlighting}{Verbatim}{commandchars=\\\{\}}
% Add ',fontsize=\small' for more characters per line
\usepackage{framed}
\definecolor{shadecolor}{RGB}{241,243,245}
\newenvironment{Shaded}{\begin{snugshade}}{\end{snugshade}}
\newcommand{\AlertTok}[1]{\textcolor[rgb]{0.68,0.00,0.00}{#1}}
\newcommand{\AnnotationTok}[1]{\textcolor[rgb]{0.37,0.37,0.37}{#1}}
\newcommand{\AttributeTok}[1]{\textcolor[rgb]{0.40,0.45,0.13}{#1}}
\newcommand{\BaseNTok}[1]{\textcolor[rgb]{0.68,0.00,0.00}{#1}}
\newcommand{\BuiltInTok}[1]{\textcolor[rgb]{0.00,0.23,0.31}{#1}}
\newcommand{\CharTok}[1]{\textcolor[rgb]{0.13,0.47,0.30}{#1}}
\newcommand{\CommentTok}[1]{\textcolor[rgb]{0.37,0.37,0.37}{#1}}
\newcommand{\CommentVarTok}[1]{\textcolor[rgb]{0.37,0.37,0.37}{\textit{#1}}}
\newcommand{\ConstantTok}[1]{\textcolor[rgb]{0.56,0.35,0.01}{#1}}
\newcommand{\ControlFlowTok}[1]{\textcolor[rgb]{0.00,0.23,0.31}{\textbf{#1}}}
\newcommand{\DataTypeTok}[1]{\textcolor[rgb]{0.68,0.00,0.00}{#1}}
\newcommand{\DecValTok}[1]{\textcolor[rgb]{0.68,0.00,0.00}{#1}}
\newcommand{\DocumentationTok}[1]{\textcolor[rgb]{0.37,0.37,0.37}{\textit{#1}}}
\newcommand{\ErrorTok}[1]{\textcolor[rgb]{0.68,0.00,0.00}{#1}}
\newcommand{\ExtensionTok}[1]{\textcolor[rgb]{0.00,0.23,0.31}{#1}}
\newcommand{\FloatTok}[1]{\textcolor[rgb]{0.68,0.00,0.00}{#1}}
\newcommand{\FunctionTok}[1]{\textcolor[rgb]{0.28,0.35,0.67}{#1}}
\newcommand{\ImportTok}[1]{\textcolor[rgb]{0.00,0.46,0.62}{#1}}
\newcommand{\InformationTok}[1]{\textcolor[rgb]{0.37,0.37,0.37}{#1}}
\newcommand{\KeywordTok}[1]{\textcolor[rgb]{0.00,0.23,0.31}{\textbf{#1}}}
\newcommand{\NormalTok}[1]{\textcolor[rgb]{0.00,0.23,0.31}{#1}}
\newcommand{\OperatorTok}[1]{\textcolor[rgb]{0.37,0.37,0.37}{#1}}
\newcommand{\OtherTok}[1]{\textcolor[rgb]{0.00,0.23,0.31}{#1}}
\newcommand{\PreprocessorTok}[1]{\textcolor[rgb]{0.68,0.00,0.00}{#1}}
\newcommand{\RegionMarkerTok}[1]{\textcolor[rgb]{0.00,0.23,0.31}{#1}}
\newcommand{\SpecialCharTok}[1]{\textcolor[rgb]{0.37,0.37,0.37}{#1}}
\newcommand{\SpecialStringTok}[1]{\textcolor[rgb]{0.13,0.47,0.30}{#1}}
\newcommand{\StringTok}[1]{\textcolor[rgb]{0.13,0.47,0.30}{#1}}
\newcommand{\VariableTok}[1]{\textcolor[rgb]{0.07,0.07,0.07}{#1}}
\newcommand{\VerbatimStringTok}[1]{\textcolor[rgb]{0.13,0.47,0.30}{#1}}
\newcommand{\WarningTok}[1]{\textcolor[rgb]{0.37,0.37,0.37}{\textit{#1}}}

\usepackage{longtable,booktabs,array}
\usepackage{calc} % for calculating minipage widths
\usepackage{caption}
% Make caption package work with longtable
\makeatletter
\def\fnum@table{\tablename~\thetable}
\makeatother
\usepackage{graphicx}
\makeatletter
\newsavebox\pandoc@box
\newcommand*\pandocbounded[1]{% scales image to fit in text height/width
  \sbox\pandoc@box{#1}%
  \Gscale@div\@tempa{\textheight}{\dimexpr\ht\pandoc@box+\dp\pandoc@box\relax}%
  \Gscale@div\@tempb{\linewidth}{\wd\pandoc@box}%
  \ifdim\@tempb\p@<\@tempa\p@\let\@tempa\@tempb\fi% select the smaller of both
  \ifdim\@tempa\p@<\p@\scalebox{\@tempa}{\usebox\pandoc@box}%
  \else\usebox{\pandoc@box}%
  \fi%
}
% Set default figure placement to htbp
\def\fps@figure{htbp}
\makeatother



\ifLuaTeX
\usepackage[bidi=basic]{babel}
\else
\usepackage[bidi=default]{babel}
\fi
% get rid of language-specific shorthands (see #6817):
\let\LanguageShortHands\languageshorthands
\def\languageshorthands#1{}


\setlength{\emergencystretch}{3em} % prevent overfull lines

\providecommand{\tightlist}{%
  \setlength{\itemsep}{0pt}\setlength{\parskip}{0pt}}



 


\setbeamertemplate{headline}{} % Elimina la cabecera del tema Madrid
\setbeamertemplate{section in head/foot}{} % Limpia la numeración de secciones
\renewcommand{\contentsname}{Índice}
\renewcommand{\refname}{Referencias}
\renewcommand{\figurename}{Figura}
\renewcommand{\tablename}{Tabla}
\logo{%
  \begin{tabular}{c}
    \includegraphics[height=1.2cm]{EscudoU.png}
  \end{tabular}%
}
% --- CONFIGURACIÓN DE COLORES INSTITUCIONALES ---
\definecolor{UnicorGreen}{HTML}{024426}
\definecolor{UnicorText}{HTML}{F9F9F9}
\definecolor{UnicorWhite}{HTML}{FFFFFF}
\definecolor{UnicorGold}{HTML}{D1B061}

% Aplicación al fondo y texto general
\setbeamercolor{background canvas}{bg=UnicorGreen}
\setbeamercolor{normal text}{fg=UnicorText}

% Aplicación a títulos y cabeceras
%\setbeamercolor{title}{fg=UnicorWhite}
\setbeamercolor{frametitle}{fg=UnicorWhite}
\setbeamercolor{framesubtitle}{fg=UnicorWhite}

% Aplicación a viñetas y estructura
\setbeamercolor{structure}{fg=UnicorWhite}
\setbeamercolor{itemize item}{fg=UnicorWhite}
\setbeamercolor{itemize subitem}{fg=UnicorWhite}
\setbeamercolor{enumerate item}{fg=UnicorWhite}

% Forzar a LaTeX a repintar el texto general al iniciar el documento
\AtBeginDocument{\usebeamercolor[fg]{normal text}}

% Aplicación a hipervínculos
\hypersetup{colorlinks=true, urlcolor=UnicorGold, linkcolor=UnicorWhite}

% --- ACTIVAR ESTILO DE BLOQUES Y SOMBRAS (ESTILO MADRID) ---
\useinnertheme[shadow=true]{rounded}

% --- COLORES DEL RECUADRO DEL TÍTULO PRINCIPAL ---
% Fondo dorado con texto verde
\setbeamercolor{title}{bg=UnicorGold, fg=UnicorGreen}
\setbeamercolor{frametitle}{fg=UnicorWhite} % Mantiene el título de las diapositivas normal
\setbeamercolor{framesubtitle}{fg=UnicorWhite}

% --- COLORES PARA TEOREMAS, DEFINICIONES Y BLOQUES ---
% Título del bloque (Fondo dorado, texto verde)
\setbeamercolor{block title}{bg=UnicorGold, fg=UnicorGreen}
% Cuerpo del bloque (Fondo blanco, texto verde)
\setbeamercolor{block body}{bg=UnicorWhite, fg=UnicorGreen}

\AtBeginDocument{
  \author{%
    \begin{tabular}{c}
      Mario A. Morales Rivera \\
      \footnotesize \texttt{mamorales@correo.unicordoba.edu.co} \\
      \footnotesize Universidad de Córdoba
    \end{tabular}
  }
  \institute{}
}
% --- LIMPIAR DIAPOSITIVAS DE SECCIÓN (ELIMINAR "SECTION X") ---
\setbeamertemplate{section page}{
  \begin{center}
    \begin{beamercolorbox}[sep=12pt,center,rounded=true,shadow=true]{title}
      \usebeamerfont{title}\insertsection\par
    \end{beamercolorbox}
  \end{center}
}
\makeatletter
\@ifpackageloaded{tcolorbox}{}{\usepackage[skins,breakable]{tcolorbox}}
\@ifpackageloaded{fontawesome5}{}{\usepackage{fontawesome5}}
\definecolor{quarto-callout-color}{HTML}{909090}
\definecolor{quarto-callout-note-color}{HTML}{0758E5}
\definecolor{quarto-callout-important-color}{HTML}{CC1914}
\definecolor{quarto-callout-warning-color}{HTML}{EB9113}
\definecolor{quarto-callout-tip-color}{HTML}{00A047}
\definecolor{quarto-callout-caution-color}{HTML}{FC5300}
\definecolor{quarto-callout-color-frame}{HTML}{acacac}
\definecolor{quarto-callout-note-color-frame}{HTML}{4582ec}
\definecolor{quarto-callout-important-color-frame}{HTML}{d9534f}
\definecolor{quarto-callout-warning-color-frame}{HTML}{f0ad4e}
\definecolor{quarto-callout-tip-color-frame}{HTML}{02b875}
\definecolor{quarto-callout-caution-color-frame}{HTML}{fd7e14}
\makeatother
\makeatletter
\@ifpackageloaded{caption}{}{\usepackage{caption}}
\AtBeginDocument{%
\ifdefined\contentsname
  \renewcommand*\contentsname{Tabla de contenidos}
\else
  \newcommand\contentsname{Tabla de contenidos}
\fi
\ifdefined\listfigurename
  \renewcommand*\listfigurename{Listado de Figuras}
\else
  \newcommand\listfigurename{Listado de Figuras}
\fi
\ifdefined\listtablename
  \renewcommand*\listtablename{Listado de Tablas}
\else
  \newcommand\listtablename{Listado de Tablas}
\fi
\ifdefined\figurename
  \renewcommand*\figurename{Figura}
\else
  \newcommand\figurename{Figura}
\fi
\ifdefined\tablename
  \renewcommand*\tablename{Tabla}
\else
  \newcommand\tablename{Tabla}
\fi
}
\@ifpackageloaded{float}{}{\usepackage{float}}
\floatstyle{ruled}
\@ifundefined{c@chapter}{\newfloat{codelisting}{h}{lop}}{\newfloat{codelisting}{h}{lop}[chapter]}
\floatname{codelisting}{Listado}
\newcommand*\listoflistings{\listof{codelisting}{Listado de Listados}}
\makeatother
\makeatletter
\makeatother
\makeatletter
\@ifpackageloaded{caption}{}{\usepackage{caption}}
\@ifpackageloaded{subcaption}{}{\usepackage{subcaption}}
\makeatother
\makeatletter
\@ifpackageloaded{tikz}{}{\usepackage{tikz}}
\makeatother
        \newcommand*\circled[1]{\tikz[baseline=(char.base)]{
          \node[shape=circle,draw,inner sep=1pt] (char) {{\scriptsize#1}};}}  
                  

\usepackage{bookmark}
\IfFileExists{xurl.sty}{\usepackage{xurl}}{} % add URL line breaks if available
\urlstyle{same}
\hypersetup{
  pdftitle={Data Frames de Pandas},
  pdfauthor={Mario Morales},
  pdflang={es},
  hidelinks,
  pdfcreator={LaTeX via pandoc}}


\title{Data Frames de Pandas}
\subtitle{Estadística con Python}
\author{Mario Morales}
\date{2026-09-21}

\begin{document}
\frame{\titlepage}

\renewcommand*\contentsname{Tabla de contenidos}
\begin{frame}[allowframebreaks]
  \frametitle{Tabla de contenidos}
  \setcounter{tocdepth}{3}
  \tableofcontents
\end{frame}
\setcounter{tocdepth}{3}
\tableofcontents
}

\section{Pandas DataFrames}\label{pandas-dataframes}

\begin{frame}[fragile]{Pandas DataFrames}
\begin{block}{Marcos de datos de Pandas}
\phantomsection\label{marcos-de-datos-de-pandas}
\index{Pandas!Marcos de datos}\index{Data Frames!de Pandas}\index{Pandas!DataFrames}

\begin{tcolorbox}[enhanced jigsaw, opacitybacktitle=0.6, coltitle=black, bottomrule=.15mm, opacityback=0, colback=white, colframe=quarto-callout-note-color-frame, rightrule=.15mm, left=2mm, toptitle=1mm, arc=.35mm, bottomtitle=1mm, titlerule=0mm, title=\textcolor{quarto-callout-note-color}{\faInfo}\hspace{0.5em}{Nota}, toprule=.15mm, colbacktitle=quarto-callout-note-color!10!white, leftrule=.75mm, breakable]

El marco de datos de Pandas o \emph{Pandas \texttt{DataFrame}} es la
estructura principal de datos en esta biblioteca de Python. Un marco de
datos de Pandas es un arreglo en forma de tabla (bidimensional) que
soporta columnas de (potencialmente) tipos heterogéneos de datos.

\end{tcolorbox}
\end{block}

\begin{block}{Marcos de datos de Pandas (II)}
\phantomsection\label{marcos-de-datos-de-pandas-ii}
\begin{itemize}
\item
  Los \texttt{DataFrame} son de tamaño mutable, es decir, se les puede
  agregar o quitar filas o columnas después de su creación;
\item
  Sus filas y columnas (conocidos como ejes, \emph{axis}) están
  etiquetados y se pueden efectuar operaciones aritméticas teniendo en
  cuenta dichas etiquetas.
\item
  Los \texttt{DataFrames} de Pandas se pueden considerar como un
  contenedor similar a un diccionario para objetos que son
  \texttt{Pandas\ Series}.
\end{itemize}
\end{block}
\end{frame}

\section{\texorpdfstring{Creación y operaciones con \texttt{DataFrames}
de
pandas}{Creación y operaciones con DataFrames de pandas}}\label{creaciuxf3n-y-operaciones-con-dataframes-de-pandas}

\begin{frame}[fragile]{Creación y operaciones con \texttt{DataFrames} de
pandas}
\begin{block}{Creación de Pandas \texttt{DataFrames}}
\phantomsection\label{creaciuxf3n-de-pandas-dataframes}
\index{Data Frames}
\index{Data Frames!creación}\index{Data Frames!operaciones con}

\begin{itemize}
\item
  A continuación se crea un \texttt{DataFrame} de Pandas.
\item
  Lo primero que hacemos es importar el paquete \emph{Pandas} mediante
  \texttt{import\ pandas\ as\ pd}.
\item
  Como se dijo en la clase anterior, no es obligatorio que se importe
  con el nombre \texttt{pd}, pero es lo usual en la comunidad de
  usuarios y útil para compatibilidad del código.
\item
  Luego creamos un \texttt{DataFrame} de nombre \texttt{df}, con tres
  columnas y siete filas.
\end{itemize}
\end{block}

\begin{block}{Creación de \texttt{DataFrames}}
\phantomsection\label{creaciuxf3n-de-dataframes}
\begin{Shaded}
\begin{Highlighting}[]
\ImportTok{import}\NormalTok{ pandas }\ImportTok{as}\NormalTok{ pd}
\ImportTok{import}\NormalTok{ numpy  }\ImportTok{as}\NormalTok{ np}
\NormalTok{df}\OperatorTok{=}\NormalTok{pd.DataFrame(\{}\StringTok{"col1"}\NormalTok{:[}\DecValTok{1}\NormalTok{,}\DecValTok{3}\NormalTok{,}\DecValTok{2}\NormalTok{,}\DecValTok{4}\NormalTok{,}\DecValTok{7}\NormalTok{,}\DecValTok{9}\NormalTok{,}\DecValTok{5}\NormalTok{],}
                \StringTok{"col2"}\NormalTok{:[}\DecValTok{9}\NormalTok{,}\DecValTok{0}\NormalTok{,}\DecValTok{2}\NormalTok{,}\DecValTok{1}\NormalTok{,}\DecValTok{7}\NormalTok{,}\DecValTok{3}\NormalTok{,}\DecValTok{4}\NormalTok{],}
                \StringTok{"col3"}\NormalTok{:[}\StringTok{"F"}\NormalTok{,}\StringTok{"M"}\NormalTok{,}\StringTok{"F"}\NormalTok{,}\StringTok{"F"}\NormalTok{,}\StringTok{"F"}\NormalTok{,}\StringTok{"M"}\NormalTok{,}\StringTok{"M"}\NormalTok{]\})}
\end{Highlighting}
\end{Shaded}

\begin{itemize}
\item
  Note que la función \texttt{DataFrame} toma como argumento un
  diccionario cuyos valores son listas, todas de la misma longitud, que
  pasan a ser las columnas del marco de datos.
\item
  Las claves del diccionario pasan a ser los nombres de las columnas.
\item
  Además de listas, los valores del diccionario también pueden ser
  arreglos unidimensionales de Numpy, diccionarios o Series de pandas.
\end{itemize}
\end{block}

\begin{block}{Creación de \texttt{DataFrames} (II)}
\phantomsection\label{creaciuxf3n-de-dataframes-ii}
Los valores del diccionario pueden ser tuplas

\begin{Shaded}
\begin{Highlighting}[numbers=left,,]
\NormalTok{df}\OperatorTok{=}\NormalTok{pd.DataFrame(\{}\StringTok{"col1"}\NormalTok{:(}\DecValTok{1}\NormalTok{,}\DecValTok{3}\NormalTok{,}\DecValTok{2}\NormalTok{,}\DecValTok{4}\NormalTok{,}\DecValTok{7}\NormalTok{,}\DecValTok{9}\NormalTok{,}\DecValTok{5}\NormalTok{),}
                \StringTok{"col2"}\NormalTok{:[}\DecValTok{9}\NormalTok{,}\DecValTok{0}\NormalTok{,}\DecValTok{2}\NormalTok{,}\DecValTok{1}\NormalTok{,}\DecValTok{7}\NormalTok{,}\DecValTok{3}\NormalTok{,}\DecValTok{4}\NormalTok{],}
            \StringTok{"col3"}\NormalTok{:[}\StringTok{"F"}\NormalTok{,}\StringTok{"M"}\NormalTok{,}\StringTok{"F"}\NormalTok{,}\StringTok{"F"}\NormalTok{,}\StringTok{"F"}\NormalTok{,}\StringTok{"M"}\NormalTok{,}\StringTok{"M"}\NormalTok{]\})}
\end{Highlighting}
\end{Shaded}
\end{block}

\begin{block}{Creación de \texttt{DataFrames} (II)}
\phantomsection\label{creaciuxf3n-de-dataframes-ii-1}
Los valores del diccionario pueden ser series

\begin{Shaded}
\begin{Highlighting}[numbers=left,,]
\NormalTok{Dic}\OperatorTok{=}\BuiltInTok{dict}\NormalTok{(}\BuiltInTok{zip}\NormalTok{(}\StringTok{"ABCDEFG"}\NormalTok{,(}\DecValTok{1}\NormalTok{,}\DecValTok{3}\NormalTok{,}\DecValTok{2}\NormalTok{,}\DecValTok{4}\NormalTok{,}\DecValTok{7}\NormalTok{,}\DecValTok{9}\NormalTok{,}\DecValTok{5}\NormalTok{)))}
\NormalTok{D}\OperatorTok{=}\NormalTok{pd.Series(Dic) }
\NormalTok{df}\OperatorTok{=}\NormalTok{pd.DataFrame(\{}\StringTok{"col1"}\NormalTok{:D,}
                \StringTok{"col2"}\NormalTok{:[}\DecValTok{9}\NormalTok{,}\DecValTok{0}\NormalTok{,}\DecValTok{2}\NormalTok{,}\DecValTok{1}\NormalTok{,}\DecValTok{7}\NormalTok{,}\DecValTok{3}\NormalTok{,}\DecValTok{4}\NormalTok{],}
            \StringTok{"col3"}\NormalTok{:[}\StringTok{"F"}\NormalTok{,}\StringTok{"M"}\NormalTok{,}\StringTok{"F"}\NormalTok{,}\StringTok{"F"}\NormalTok{,}\StringTok{"F"}\NormalTok{,}\StringTok{"M"}\NormalTok{,}\StringTok{"M"}\NormalTok{]\})}
\NormalTok{df}
\end{Highlighting}
\end{Shaded}

\begin{verbatim}
   col1  col2 col3
A     1     9    F
B     3     0    M
C     2     2    F
D     4     1    F
E     7     7    F
F     9     3    M
G     5     4    M
\end{verbatim}
\end{block}

\begin{block}{Creación de \texttt{DataFrames} desde diccionarios}
\phantomsection\label{creaciuxf3n-de-dataframes-desde-diccionarios}
Otra forma de crear marcos de datos de Pandas es desde una lista de
diccionarios, veamos el siguiente ejemplo:

\phantomsection\label{annotated-cell-4}%
\begin{Shaded}
\begin{Highlighting}[]
\NormalTok{d}\OperatorTok{=}\NormalTok{[\{}\StringTok{\textquotesingle{}x\textquotesingle{}}\NormalTok{:}\DecValTok{2}\NormalTok{,}\StringTok{\textquotesingle{}y\textquotesingle{}}\NormalTok{:}\DecValTok{0}\NormalTok{,}\StringTok{\textquotesingle{}z\textquotesingle{}}\NormalTok{:}\DecValTok{1}\NormalTok{\},\{}\StringTok{\textquotesingle{}x\textquotesingle{}}\NormalTok{:}\DecValTok{3}\NormalTok{,}\StringTok{\textquotesingle{}y\textquotesingle{}}\NormalTok{:}\DecValTok{1}\NormalTok{,}\StringTok{\textquotesingle{}z\textquotesingle{}}\NormalTok{:}\DecValTok{2}\NormalTok{\},\{}\StringTok{\textquotesingle{}x\textquotesingle{}}\NormalTok{:}\DecValTok{6}\NormalTok{,}\StringTok{\textquotesingle{}y\textquotesingle{}}\NormalTok{:}\DecValTok{4}\NormalTok{\}] }\hspace*{\fill}\NormalTok{\circled{1}}
\NormalTok{pd.DataFrame(d) }\hspace*{\fill}\NormalTok{\circled{2}}
\end{Highlighting}
\end{Shaded}

\begin{description}
\tightlist
\item[\circled{1}]
Lista de diccionarios, cada diccionario es una fila del marco de datos.
Las claves pasan a ser los nombres de columnas.
\item[\circled{2}]
Se crea el marco de datos.
\end{description}

\begin{verbatim}
   x  y    z
0  2  0  1.0
1  3  1  2.0
2  6  4  NaN
\end{verbatim}
\end{block}

\begin{block}{Creación de \texttt{DataFrames} desde arreglos de NumPy}
\phantomsection\label{creaciuxf3n-de-dataframes-desde-arreglos-de-numpy}
También podemos crear marcos de datos de Pandas a partir de arreglos
bidimencionales de Numpy, como se muestra a continuación

\phantomsection\label{annotated-cell-5}%
\begin{Shaded}
\begin{Highlighting}[]
\NormalTok{ar}\OperatorTok{=}\NormalTok{np.array([[}\DecValTok{2}\NormalTok{,}\DecValTok{0}\NormalTok{,}\DecValTok{1}\NormalTok{],[}\DecValTok{3}\NormalTok{,}\DecValTok{1}\NormalTok{,}\DecValTok{2}\NormalTok{],[}\DecValTok{6}\NormalTok{,}\DecValTok{4}\NormalTok{,}\DecValTok{1}\NormalTok{],[}\DecValTok{4}\NormalTok{,}\DecValTok{1}\NormalTok{,}\DecValTok{2}\NormalTok{]]) }\hspace*{\fill}\NormalTok{\circled{1}}
\NormalTok{ardf}\OperatorTok{=}\NormalTok{pd.DataFrame(ar) }\hspace*{\fill}\NormalTok{\circled{2}}
\NormalTok{ardf.columns}\OperatorTok{=}\BuiltInTok{list}\NormalTok{([}\StringTok{"A"}\NormalTok{,}\StringTok{"B"}\NormalTok{,}\StringTok{"C"}\NormalTok{]) }\hspace*{\fill}\NormalTok{\circled{3}}
\end{Highlighting}
\end{Shaded}

\begin{description}
\tightlist
\item[\circled{1}]
Arreglo de 4 filas 3 columnas.
\item[\circled{2}]
Creamos un \texttt{Data\ Frame} de Pandas.
\item[\circled{3}]
Asignamos nombres a las columnas del \texttt{Data\ Frame}.
\end{description}
\end{block}

\begin{block}{Creación de \texttt{DataFrames} desde arreglos de NumPy}
\phantomsection\label{creaciuxf3n-de-dataframes-desde-arreglos-de-numpy-1}
Los nombres de las columnas también se pueden asignar al momento de
crear el marco de datos, mediante el argumento
\texttt{columns={[}"A","B","C"{]}}.

\phantomsection\label{annotated-cell-6}%
\begin{Shaded}
\begin{Highlighting}[]
\NormalTok{ar}\OperatorTok{=}\NormalTok{np.array([[}\DecValTok{2}\NormalTok{,}\DecValTok{0}\NormalTok{,}\DecValTok{1}\NormalTok{],[}\DecValTok{3}\NormalTok{,}\DecValTok{1}\NormalTok{,}\DecValTok{2}\NormalTok{],[}\DecValTok{6}\NormalTok{,}\DecValTok{4}\NormalTok{,}\DecValTok{1}\NormalTok{],[}\DecValTok{4}\NormalTok{,}\DecValTok{1}\NormalTok{,}\DecValTok{2}\NormalTok{]]) }\hspace*{\fill}\NormalTok{\circled{1}}
\NormalTok{ardf}\OperatorTok{=}\NormalTok{pd.DataFrame(ar,columns}\OperatorTok{=}\NormalTok{[}\StringTok{"A"}\NormalTok{,}\StringTok{"B"}\NormalTok{,}\StringTok{"C"}\NormalTok{]) }\hspace*{\fill}\NormalTok{\circled{2}}
\end{Highlighting}
\end{Shaded}

\begin{description}
\tightlist
\item[\circled{1}]
Arreglo de 4 filas 3 columnas.
\item[\circled{2}]
Creamos un \texttt{Data\ Frame} de Pandas asignando nombres de columnas.
\end{description}
\end{block}

\begin{block}{Creación de \texttt{DataFrames} (IV)}
\phantomsection\label{creaciuxf3n-de-dataframes-iv}
\begin{itemize}
\item
  \texttt{DataFrame.from\_dict}, crea un marco de datos de Pandas a
  partir de un diccionario,
\item
  Contamos con el argumento \texttt{orient}

  \begin{itemize}
  \tightlist
  \item
    ``columns'' por defecto, usa las claves como etiquetas columnas
  \item
    ``index'' para usar las claves del diccionario como etiquetas de
    fila.
  \end{itemize}
\end{itemize}
\end{block}
\end{frame}

\begin{frame}[fragile]{}
\phantomsection\label{section}
\phantomsection\label{annotated-cell-7}%
\begin{Shaded}
\begin{Highlighting}[]
\NormalTok{d}\OperatorTok{=}\NormalTok{\{}\StringTok{\textquotesingle{}A\textquotesingle{}}\NormalTok{: [}\DecValTok{1}\NormalTok{, }\DecValTok{2}\NormalTok{, }\DecValTok{3}\NormalTok{], }\StringTok{\textquotesingle{}B\textquotesingle{}}\NormalTok{: [}\DecValTok{4}\NormalTok{, }\DecValTok{5}\NormalTok{, }\DecValTok{6}\NormalTok{]\} }\hspace*{\fill}\NormalTok{\circled{1}}
\NormalTok{dfd}\OperatorTok{=}\NormalTok{pd.DataFrame.from\_dict(d,orient}\OperatorTok{=}\StringTok{\textquotesingle{}index\textquotesingle{}}\NormalTok{) }\hspace*{\fill}\NormalTok{\circled{2}}
\NormalTok{dfd.columns}\OperatorTok{=}\NormalTok{[}\StringTok{"uno"}\NormalTok{,}\StringTok{"dos"}\NormalTok{,}\StringTok{"tres"}\NormalTok{] }\hspace*{\fill}\NormalTok{\circled{3}}
\end{Highlighting}
\end{Shaded}

\begin{description}
\tightlist
\item[\circled{1}]
Diccionario con listas como valores.
\item[\circled{2}]
La opción \texttt{orient=\textquotesingle{}index\textquotesingle{}} hace
que los valores del diccionario, que son las listas, pasen a ser las
filas del marco de datos.
\item[\circled{3}]
Asignamos nombres a las columnas. También podemos hacerlo mediante,
\texttt{columns={[}"uno","dos","tres"{]}}, dentro del constructor
\texttt{from\_dict}.
\end{description}
\end{frame}

\section{\texorpdfstring{Inspeccionando
\texttt{Data\ Frames}}{Inspeccionando Data Frames}}\label{inspeccionando-data-frames}

\begin{frame}[fragile]{Inspeccionando \texttt{Data\ Frames}}
\begin{block}{Inspeccionando \texttt{Data\ Frames}}
\phantomsection\label{inspeccionando-data-frames-1}
Con los métodos \texttt{.head(k)} y \texttt{.tail(k)} podemos examinar
las \(k\) primeras y últimas filas del \texttt{DataFrame},
respectivamente. \index{Data Frames!Método!head}
\index{Data Frames!Método!tail}

\begin{Shaded}
\begin{Highlighting}[]
\NormalTok{df.head(}\DecValTok{4}\NormalTok{) }\CommentTok{\# 4 primeras filas }
\end{Highlighting}
\end{Shaded}

\begin{verbatim}
   col1  col2 col3
A     1     9    F
B     3     0    M
C     2     2    F
D     4     1    F
\end{verbatim}

\begin{Shaded}
\begin{Highlighting}[]
\NormalTok{df.tail(}\DecValTok{3}\NormalTok{) }\CommentTok{\# 3 últimas filas }
\end{Highlighting}
\end{Shaded}

\begin{verbatim}
   col1  col2 col3
E     7     7    F
F     9     3    M
G     5     4    M
\end{verbatim}
\end{block}

\begin{block}{Método \texttt{.info}}
\phantomsection\label{muxe9todo-.info}
Proporciona detalles esenciales como:

\begin{itemize}
\item
  La clase de objeto, el número de filas y columnas,
\item
  Cuántos datos disponibles tiene la columna,
\item
  El tipo de datos que contiene cada columna y la cantidad de memoria
  ocupada por el marco de datos.
\end{itemize}

\begin{Shaded}
\begin{Highlighting}[]
\NormalTok{df.info()}
\end{Highlighting}
\end{Shaded}

\begin{verbatim}
<class 'pandas.core.frame.DataFrame'>
Index: 7 entries, A to G
Data columns (total 3 columns):
 #   Column  Non-Null Count  Dtype 
---  ------  --------------  ----- 
 0   col1    7 non-null      int64 
 1   col2    7 non-null      int64 
 2   col3    7 non-null      object
dtypes: int64(2), object(1)
memory usage: 224.0+ bytes
\end{verbatim}
\end{block}

\begin{block}{Método \texttt{.shape}}
\phantomsection\label{muxe9todo-.shape}
Recupera el número de filas y columnas de la tabla de datos

\begin{itemize}
\item
  Regresa una tupla con dos elementos,

  \begin{itemize}
  \tightlist
  \item
    El primero corresponde al número de filas y
  \item
    El segundo al número de columnas
  \end{itemize}
\end{itemize}

\phantomsection\label{annotated-cell-11}%
\begin{Shaded}
\begin{Highlighting}[]
\NormalTok{nfil,ncol }\OperatorTok{=}\NormalTok{ df.shape }\hspace*{\fill}\NormalTok{\circled{1}}
\NormalTok{df.shape }\hspace*{\fill}\NormalTok{\circled{2}}
\end{Highlighting}
\end{Shaded}

\begin{description}
\tightlist
\item[\circled{1}]
En \texttt{nfil} tenemos el número de filas y en \texttt{ncol} el número
de columnas de \texttt{df}.\\
\item[\circled{2}]
Regresa una tupla con dos elementos, el primero contiene el numero de
filas y el segundo el numero de columnas de \texttt{df}.
\end{description}
\end{block}

\begin{block}{Número de filas y columnas}
\phantomsection\label{nuxfamero-de-filas-y-columnas}
\begin{itemize}
\item
  Podemos guardar esa información en variables de Python para su uso en
  alguna rutina, mediante \texttt{nfil,ncol\ =\ df.shape}, que es justo
  lo que se hizo en la primera línea del bloque de código anterior.
\item
  Se observa, en la salida, que el marco de datos \texttt{df} tiene 7 y
  3 columnas.
\end{itemize}
\end{block}

\begin{block}{Método \texttt{.index}}
\phantomsection\label{muxe9todo-.index}
Regresa un objeto de clase \texttt{RangeIndex} que contiene la
información de los índices asociados a las filas del marco de datos. Con
la función \texttt{list()} podemos obtener una lista que contiene los
índices

\begin{Shaded}
\begin{Highlighting}[]
\BuiltInTok{list}\NormalTok{(dfd.index) }
\end{Highlighting}
\end{Shaded}

\begin{verbatim}
['A', 'B']
\end{verbatim}
\end{block}

\begin{block}{El método \texttt{.columns}}
\phantomsection\label{el-muxe9todo-.columns}
Permite recuperar y cambiar los nombres de las columnas del marco de
datos.

\begin{Shaded}
\begin{Highlighting}[]
\NormalTok{dfd.columns }
\end{Highlighting}
\end{Shaded}

\begin{verbatim}
Index(['uno', 'dos', 'tres'], dtype='object')
\end{verbatim}

\begin{Shaded}
\begin{Highlighting}[]
\BuiltInTok{list}\NormalTok{(dfd.columns) }
\end{Highlighting}
\end{Shaded}

\begin{verbatim}
['uno', 'dos', 'tres']
\end{verbatim}
\end{block}
\end{frame}

\section{Creación a partir de fuentes
externas}\label{creaciuxf3n-a-partir-de-fuentes-externas}

\begin{frame}[fragile]{Creación a partir de fuentes externas}
\begin{block}{Lectura de datos desde archivos externos.}
\phantomsection\label{sec-lecturaDatosExternos}
\index{Lectura!de datos externos} Es posible crear
\texttt{Pandas\ DataFrames} a partir de datos que se encuentran en
fuentes externas, tales como archivos de \texttt{MS\ Excel}
(\texttt{.xls} o \texttt{.xlsx}), Open Document Format (\texttt{.odf}),
ficheros de texto plano (\texttt{.txt} o \texttt{.csv}), en bases de
datos relacionales \texttt{SQL} e incluso en tablas de páginas web. La
lectura de datos en el formato correcto y verificar que no hayan errores
en el proceso, es la parte que toma más tiempo en el análisis de datos.
\end{block}

\begin{block}{Lectura de datos desde archivos externos (II)}
\phantomsection\label{lectura-de-datos-desde-archivos-externos-ii}
Cuando los datos están organizados en un archivo de texto plano,
deberíamos iniciar siempre con una inspección visual de los datos, en
particular deberíamos chequear:

\begin{itemize}
\item
  ¿Tienen los datos un encabezado al inicio o texto al final?
\item
  Lo usual es que el encabezado tenga los nombres de las columnas,
\item
  Las notas al final de los archivos se usan para dar una descripción de
  los datos.
\end{itemize}

Por ejemplo, se explica el contenido de las columnas, si hay alguna
codificación de las respuestas que haya que tener en cuenta, entre otras
cosas.
\end{block}

\begin{block}{Lectura de datos desde archivos externos (III)}
\phantomsection\label{lectura-de-datos-desde-archivos-externos-iii}
\begin{itemize}
\tightlist
\item
  ¿Hay líneas vacías al final del archivo?
\item
  ¿Hay espacios en blanco antes del primer número o al final de cada
  línea?
\item
  ¿Están los datos separados por espacios, tabuladores, comas (,) o
  punto y coma (;)?
\item
  ¿Qué carácter o símbolo se está usando como separador de decimales?
\item
  ¿Qué carácter o símbolo se está usando para agrupar cadenas de texto?
\end{itemize}
\end{block}

\begin{block}{Lectura de datos desde archivos externos (IV)}
\phantomsection\label{lectura-de-datos-desde-archivos-externos-iv}
La lectura con Pandas tiene los siguientes pasos:

\begin{itemize}
\item
  Cambiar el directorio actual (carpeta de trabajo) a la ruta donde está
  el archivo que contiene los datos.
\item
  Listar los archivos de la carpeta de trabajo.
\item
  Este paso no es obligatorio, pero es útil para verificar que el
  archivo se encuentra en esa ruta.
\item
  Leer los datos.
\item
  Verificar que los datos han sido leídos completamente y en el formato
  correcto.
\end{itemize}
\end{block}

\begin{block}{Lectura de datos desde archivos de texto plano.}
\phantomsection\label{lectura-de-datos-desde-archivos-de-texto-plano.}
\begin{itemize}
\item
  Seguidamente se importaran a un \texttt{DataFrame} de Pandas los datos
  contenidos en el archivo \texttt{DatosEst.csv}, el cual puede
  descargar del repositorio de archivos asociados al libro.
\item
  El archivo contiene información recolectada mediante una encuesta
  aplicada a estudiantes de primer semestre de la carrera de Estadística
  de la Universidad de Córdoba.
\item
  La columnas contienen la edad, el género, la estatura, el peso, una
  valoración del nivel de satisfacción del estudiante en relación al
  servicio prestado por la Universidad, entre otras características.
\end{itemize}
\end{block}

\begin{block}{Lectura de datos desde archivos de texto plano.}
\phantomsection\label{lectura-de-datos-desde-archivos-de-texto-plano.-1}
Lo primero que hacemos es examinar el directorio de trabajo actual, la
carpeta desde donde Pandas intentará leer los datos.

\phantomsection\label{annotated-cell-15}%
\begin{Shaded}
\begin{Highlighting}[]
\ImportTok{import}\NormalTok{ os }\hspace*{\fill}\NormalTok{\circled{1}}
\NormalTok{dir\_act}\OperatorTok{=}\NormalTok{os.path.abspath(os.curdir) }\hspace*{\fill}\NormalTok{\circled{2}}
\NormalTok{wd}\OperatorTok{=}\NormalTok{os.getcwd() }\hspace*{\fill}\NormalTok{\circled{3}}
\end{Highlighting}
\end{Shaded}

\begin{description}
\tightlist
\item[\circled{1}]
Se carga el modulo \texttt{os}, con el cual interactuamos, desde python,
con el sistema operativo.
\item[\circled{2}]
Obtenemos la ruta absoluta del directorio actual.
\item[\circled{3}]
Equivalente a la línea anterior.
\end{description}
\end{block}

\begin{block}{}
\phantomsection\label{section-1}
La salida del comando anterior es la ruta completa al directorio de
trabajo actual, algo como \texttt{/home/proyectos/python}, dependiendo
de su sistema operativo y la forma como tenga organizadas sus carpetas.
En caso que se quiera cambiar la carpeta de trabajo se usa la función
\texttt{chdir}, por ejemplo, para cambiar a la carpeta
\texttt{/home/proyectos/datos} se usa la siguiente línea de código

\begin{Shaded}
\begin{Highlighting}[]
\NormalTok{os.chdir(}\StringTok{\textquotesingle{}/home/proyectos/datos\textquotesingle{}}\NormalTok{) }
\end{Highlighting}
\end{Shaded}
\end{block}

\begin{block}{Lista de archivos en el directorio}
\phantomsection\label{lista-de-archivos-en-el-directorio}
Podemos obtener un listado de los archivos contenidos en la carpeta de
trabajo actual mediante

\begin{Shaded}
\begin{Highlighting}[]
\NormalTok{os.listdir()}
\end{Highlighting}
\end{Shaded}

\begin{verbatim}
['Clase4Diccionarios.pdf', 'EscudoU.png', 'clasepythoncomocalculadora.ipynb', 'DatosEst.csv']
\end{verbatim}
\end{block}

\begin{block}{Verificar si el archivo está en el directorio actual}
\phantomsection\label{verificar-si-el-archivo-estuxe1-en-el-directorio-actual}
En este punto podemos verificar si el archivo de datos que pretendemos
leer se encuentra en el listado anterior. Si el número de archivos
contenido en la carpeta es muy grande, podría ser útil la siguiente
línea, que regresaría \texttt{True} si el archivo \texttt{DatosEst.csv}
se encuentra en la carpeta y \texttt{False} en caso contrario.

\begin{Shaded}
\begin{Highlighting}[]
\CommentTok{\textquotesingle{}DatosEst.csv\textquotesingle{}} \KeywordTok{in}\NormalTok{ os.listdir()}
\end{Highlighting}
\end{Shaded}

\begin{verbatim}
True
\end{verbatim}
\end{block}

\begin{block}{Lectura del archivo}
\phantomsection\label{lectura-del-archivo}
Luego de fijar la carpeta de trabajo y estando seguros que el archivo de
datos a leer se encuentra en ella, procedemos a la lectura de los datos

\begin{Shaded}
\begin{Highlighting}[]
\NormalTok{df}\OperatorTok{=}\NormalTok{pd.read\_csv(}\StringTok{\textquotesingle{}DatosEst.csv\textquotesingle{}}\NormalTok{)}
\end{Highlighting}
\end{Shaded}
\end{block}

\begin{block}{Argumento \texttt{header}}
\phantomsection\label{argumento-header}
La función \texttt{read\_csv} está configurada, por defecto, para leer
archivos que tienen, en la primera fila, los nombres de las columnas. En
caso que eso no sea así, tendríamos que indicarlo a la función mediante
el argumento \texttt{header}, de la siguiente forma

\begin{Shaded}
\begin{Highlighting}[]
\NormalTok{df}\OperatorTok{=}\NormalTok{pd.read\_csv(}\StringTok{\textquotesingle{}DatosEst.csv\textquotesingle{}}\NormalTok{,header}\OperatorTok{=}\VariableTok{None}\NormalTok{)}
\end{Highlighting}
\end{Shaded}
\end{block}

\begin{block}{Argumento \texttt{delimiter}}
\phantomsection\label{argumento-delimiter}
Con el argumento \texttt{delimiter} se establece el delimitador de
campos, es decir, el carácter o símbolo que separa las columnas y viene,
de manera predeterminada, configurada con la coma \texttt{(,)}. En caso
que el archivo de datos a leer tenga otro delimitador, como por ejemplo
punto y coma \texttt{(;)}, debemos usar
\texttt{delimiter=\textquotesingle{};\textquotesingle{}}, como se
muestra a continuación

\begin{Shaded}
\begin{Highlighting}[]
\NormalTok{df\_es}\OperatorTok{=}\NormalTok{pd.read\_csv(}\StringTok{\textquotesingle{}DatosEst.csv\textquotesingle{}}\NormalTok{,delimiter}\OperatorTok{=}\StringTok{\textquotesingle{};\textquotesingle{}}\NormalTok{)}
\end{Highlighting}
\end{Shaded}

Para más información de la opciones de \texttt{read\_csv}, consultar su
documentación.
\end{block}

\begin{block}{Lectura de datos contenidos en archivos de excel (.xls)}
\phantomsection\label{lectura-de-datos-contenidos-en-archivos-de-excel-.xls}
Siempre es posible guardar un archivo de Microsoft Excel en formato
\texttt{csv} o \texttt{txt}. Así que, una forma de leer datos que vienen
en formato \texttt{xls} o \texttt{xlsx} es exportándolos a texto plano y
luego aplicar lo aprendido en la sección anterior.

Si por alguna razón no es posible la conversión a texto plano, en Pandas
contamos con la clase \texttt{ExelFile}, que se recomienda para leer un
libro con diferentes formatos de datos en las hojas; y la función
\texttt{read\_excel}, recomendada para leer libros con formatos de datos
idénticos en las hojas.
\end{block}

\begin{block}{Lectura de datos contenidos en archivos de excel (.xls)}
\phantomsection\label{lectura-de-datos-contenidos-en-archivos-de-excel-.xls-1}
\begin{itemize}
\item
  Suponga que queremos importar los datos que se encuentran en la
  \texttt{hoja1} del archivo \texttt{datosExel.xls}.
\item
  Luego de colocar el archivo en la carpeta de trabajo actual,
  procedemos a leerlo con la clase \texttt{ExelFile} de Pandas
\end{itemize}

\begin{Shaded}
\begin{Highlighting}[]
\NormalTok{xls}\OperatorTok{=}\NormalTok{pd.ExcelFile(}\StringTok{\textquotesingle{}datosExel.xls\textquotesingle{}}\NormalTok{)}
\NormalTok{df\_ex}\OperatorTok{=}\NormalTok{xls.parse(}\StringTok{\textquotesingle{}hoja1\textquotesingle{}}\NormalTok{,na\_values}\OperatorTok{=}\StringTok{\textquotesingle{}NA\textquotesingle{}}\NormalTok{)}
\end{Highlighting}
\end{Shaded}
\end{block}

\begin{block}{Lectura de datos contenidos en archivos de excel (.xls)}
\phantomsection\label{lectura-de-datos-contenidos-en-archivos-de-excel-.xls-2}
Ahora, a continuación, usaremos la función \texttt{read\_excel} para
leer los datos contenidos en el mismo archivo anterior

\begin{Shaded}
\begin{Highlighting}[]
\NormalTok{df\_ex}\OperatorTok{=}\NormalTok{pd.read\_excel(}\StringTok{\textquotesingle{}datosExel.xls\textquotesingle{}}\NormalTok{,}\StringTok{\textquotesingle{}hoja1\textquotesingle{}}\NormalTok{,na\_values}\OperatorTok{=}\NormalTok{[}\StringTok{\textquotesingle{}NA\textquotesingle{}}\NormalTok{])}
\NormalTok{df\_ex.head() }
\end{Highlighting}
\end{Shaded}
\end{block}
\end{frame}

\section{\texorpdfstring{Operaciones con
\texttt{DataFrames}}{Operaciones con DataFrames}}\label{operaciones-con-dataframes}

\begin{frame}[fragile]{Operaciones con \texttt{DataFrames}}
\begin{block}{Selección, adición y borrado de columnas.}
\phantomsection\label{selecciuxf3n-adiciuxf3n-y-borrado-de-columnas.}
\begin{itemize}
\item
  Con \texttt{df{[}\textquotesingle{}x\textquotesingle{}{]}} podemos
  seleccionar la columna de nombre \texttt{x}, del
  \texttt{Pandas\ DataFrame}, \texttt{df}.
\item
  De manera equivalente, mediante \texttt{df.x}.
\item
  A continuación seleccionamos \texttt{col1} del marco de datos
  \texttt{df}
\end{itemize}

\begin{Shaded}
\begin{Highlighting}[]
\NormalTok{df[}\StringTok{\textquotesingle{}col1\textquotesingle{}}\NormalTok{] }
\end{Highlighting}
\end{Shaded}

\begin{verbatim}
A    1
B    3
C    2
D    4
E    7
F    9
G    5
Name: col1, dtype: int64
\end{verbatim}
\end{block}

\begin{block}{Selección, adición y borrado de columnas (II)}
\phantomsection\label{selecciuxf3n-adiciuxf3n-y-borrado-de-columnas-ii}
\begin{itemize}
\tightlist
\item
  Observe que la clase del objeto recuperado es una serie de Pandas.
\end{itemize}

\begin{Shaded}
\begin{Highlighting}[]
\BuiltInTok{type}\NormalTok{(df[}\StringTok{\textquotesingle{}col1\textquotesingle{}}\NormalTok{])}
\end{Highlighting}
\end{Shaded}

\begin{verbatim}
<class 'pandas.core.series.Series'>
\end{verbatim}

\begin{itemize}
\tightlist
\item
  Otra forma de conseguir lo mismo es mediante \texttt{df.col1}.
\end{itemize}
\end{block}

\begin{block}{Selección de mas de una columna}
\phantomsection\label{selecciuxf3n-de-mas-de-una-columna}
\begin{itemize}
\tightlist
\item
  Si necesita extraer más de una columna al tiempo, por ejemplo, se
  puede extraer de \texttt{df} las columnas de nombre \texttt{col1} y
  \texttt{col2}, mediante
  \texttt{df{[}{[}\textquotesingle{}col1\textquotesingle{},\textquotesingle{}col2\textquotesingle{}{]}{]}},
  veamos
\end{itemize}

\begin{Shaded}
\begin{Highlighting}[]
\NormalTok{df[[}\StringTok{\textquotesingle{}col1\textquotesingle{}}\NormalTok{,}\StringTok{\textquotesingle{}col1\textquotesingle{}}\NormalTok{]].head(}\DecValTok{3}\NormalTok{) }
\end{Highlighting}
\end{Shaded}

\begin{verbatim}
   col1  col1
A     1     1
B     3     3
C     2     2
\end{verbatim}
\end{block}

\begin{block}{Adición de columnas}
\phantomsection\label{adiciuxf3n-de-columnas}
A continuación se crea la columna \texttt{rcol2} calculando la raíz
cuadrada positiva de la columna \texttt{col2}, es decir, se aplica la
transformación \(y=\sqrt x\), luego se agrega, como una columna, al
\texttt{DataFrame} \texttt{df}

\begin{Shaded}
\begin{Highlighting}[]
\NormalTok{df[}\StringTok{\textquotesingle{}rcol2\textquotesingle{}}\NormalTok{]}\OperatorTok{=}\NormalTok{np.sqrt(df[}\StringTok{"col2"}\NormalTok{])}
\NormalTok{df.head(}\DecValTok{3}\NormalTok{)}
\end{Highlighting}
\end{Shaded}

\begin{verbatim}
   col1  col2 col3     rcol2
A     1     9    F  3.000000
B     3     0    M  0.000000
C     2     2    F  1.414214
\end{verbatim}

note que ahora \texttt{df} tiene cuatro columnas.
\end{block}

\begin{block}{Adición de columnas (II)}
\phantomsection\label{adiciuxf3n-de-columnas-ii}
Otra forma de llevar a cabo la tarea anterior es mediante

\begin{Shaded}
\begin{Highlighting}[]
\NormalTok{df[}\StringTok{\textquotesingle{}rcol1\textquotesingle{}}\NormalTok{]}\OperatorTok{=}\NormalTok{np.sqrt(df.col2)  }
\end{Highlighting}
\end{Shaded}
\end{block}

\begin{block}{Adición de columnas (III)}
\phantomsection\label{adiciuxf3n-de-columnas-iii}
Seguidamente adicionamos la columna \texttt{sum\_col} al marco de datos
\texttt{df}, \texttt{sum\_col} se obtiene como la suma de las columnas
\texttt{col1} y \texttt{col2}

\begin{Shaded}
\begin{Highlighting}[]
\NormalTok{df[}\StringTok{\textquotesingle{}sum\_col\textquotesingle{}}\NormalTok{]}\OperatorTok{=}\NormalTok{df.col1 }\OperatorTok{+}\NormalTok{ df.col2 }
\NormalTok{df.head(}\DecValTok{3}\NormalTok{)}
\end{Highlighting}
\end{Shaded}

\begin{verbatim}
   col1  col2 col3     rcol2  sum_col
A     1     9    F  3.000000       10
B     3     0    M  0.000000        3
C     2     2    F  1.414214        4
\end{verbatim}
\end{block}

\begin{block}{Adición de columnas (IV)}
\phantomsection\label{adiciuxf3n-de-columnas-iv}
Como un ejemplo adicional, podríamos agregar una columna indicadora con
valor cero, en caso que \texttt{col3} sea igual a \texttt{F} y uno en
caso que sea igual a \texttt{M}. En esta oportunidad usaremos el
concepto de comprensión de listas, visto en la sección

\begin{Shaded}
\begin{Highlighting}[]
\NormalTok{df[}\StringTok{\textquotesingle{}Ind\textquotesingle{}}\NormalTok{]}\OperatorTok{=}\NormalTok{[}\DecValTok{0} \ControlFlowTok{if}\NormalTok{ x  }\ControlFlowTok{else} \DecValTok{1} \ControlFlowTok{for}\NormalTok{ x }\KeywordTok{in}\NormalTok{ df[}\StringTok{"col3"}\NormalTok{]}\OperatorTok{==}\StringTok{"F"}\NormalTok{]}
\NormalTok{df}
\end{Highlighting}
\end{Shaded}

\begin{verbatim}
   col1  col2 col3     rcol2  sum_col  Ind
A     1     9    F  3.000000       10    0
B     3     0    M  0.000000        3    1
C     2     2    F  1.414214        4    0
D     4     1    F  1.000000        5    0
E     7     7    F  2.645751       14    0
F     9     3    M  1.732051       12    1
G     5     4    M  2.000000        9    1
\end{verbatim}
\end{block}

\begin{block}{Borrado de columnas}
\phantomsection\label{borrado-de-columnas}
Para borrar una columna se usa la función \texttt{del}, como se muestra
seguidamente

\begin{Shaded}
\begin{Highlighting}[]
\KeywordTok{del}\NormalTok{ df[}\StringTok{\textquotesingle{}Ind\textquotesingle{}}\NormalTok{]}
\NormalTok{df }
\end{Highlighting}
\end{Shaded}

\begin{verbatim}
   col1  col2 col3     rcol2  sum_col
A     1     9    F  3.000000       10
B     3     0    M  0.000000        3
C     2     2    F  1.414214        4
D     4     1    F  1.000000        5
E     7     7    F  2.645751       14
F     9     3    M  1.732051       12
G     5     4    M  2.000000        9
\end{verbatim}
\end{block}
\end{frame}

\section{Filtrado de datos con
pandas}\label{filtrado-de-datos-con-pandas}

\begin{frame}[fragile]{Filtrado de datos con pandas}
\begin{block}{Selección de filas y columnas con \texttt{iloc} y
\texttt{loc}}
\phantomsection\label{selecciuxf3n-de-filas-y-columnas-con-iloc-y-loc}
\begin{itemize}
\item
  El método \texttt{iloc} selecciona tanto filas como columnas basado en
  la posición del elemento.
\item
  El formato de la sintaxis es:
  \texttt{df.iloc{[}\textless{}fila\textgreater{},\textless{}columna\textgreater{}{]}}
\item
  \texttt{fila} y/o \texttt{columna} corresponde a un número entero o
  una lista de enteros que indican la posición (índice) de las filas o
  columnas a seleccionar.
\item
  Esta forma de seleccionar filas con \texttt{.iloc} también soporta
  índices negativos, en cuyo caso la selección se realiza empezando por
  la última fila (o columna).
\item
  Cuando se entrega un solo valor, por ejemplo
  \texttt{df.iloc{[}valor{]}} la selección se hace sobre las filas,
  veamos su funcionamiento con ejemplos concretos.
\end{itemize}
\end{block}

\begin{block}{Selección de filas y columnas con \texttt{iloc} y
\texttt{loc} (II)}
\phantomsection\label{selecciuxf3n-de-filas-y-columnas-con-iloc-y-loc-ii}
\phantomsection\label{annotated-cell-32}%
\begin{Shaded}
\begin{Highlighting}[]
\NormalTok{df.iloc[}\DecValTok{0}\NormalTok{:}\DecValTok{3}\NormalTok{] }\hspace*{\fill}\NormalTok{\circled{1}}
\NormalTok{df.iloc[}\DecValTok{0}\NormalTok{] }\hspace*{\fill}\NormalTok{\circled{2}}
\NormalTok{df.iloc[[}\DecValTok{0}\NormalTok{,}\DecValTok{4}\NormalTok{,}\DecValTok{6}\NormalTok{]] }\hspace*{\fill}\NormalTok{\circled{3}}
\NormalTok{df.iloc[}\OperatorTok{{-}}\DecValTok{1}\NormalTok{] }\hspace*{\fill}\NormalTok{\circled{4}}
\end{Highlighting}
\end{Shaded}

\begin{description}
\tightlist
\item[\circled{1}]
Se seleccionan las tres primeras filas de \texttt{df}.
\item[\circled{2}]
Se selecciona la primera fila de \texttt{df}, esto es equivalente a
\texttt{df.iloc{[}0,{]}}.
\item[\circled{3}]
Selección de filas con índices específicos; se seleccionan las filas de
índice 0, 4 y 6.
\item[\circled{4}]
Seleccionamos la última fila; con \texttt{df.iloc{[}-2{]}} se selecciona
la penúltima y así sucesivamente.
\end{description}
\end{block}

\begin{block}{Selección de columnas con \texttt{.iloc}}
\phantomsection\label{selecciuxf3n-de-columnas-con-.iloc}
Para la selección de columnas se procede de forma análoga. Por ejemplo,
a continuación seleccionamos las dos primeras columnas, los dos puntos
(\texttt{:}) antes de la coma, hacen referencia a \emph{todas las
filas}.

\begin{Shaded}
\begin{Highlighting}[]
\NormalTok{df.iloc[:,}\DecValTok{0}\NormalTok{:}\DecValTok{2}\NormalTok{]}
\end{Highlighting}
\end{Shaded}

\begin{verbatim}
   col1  col2
A     1     9
B     3     0
C     2     2
D     4     1
E     7     7
F     9     3
G     5     4
\end{verbatim}
\end{block}

\begin{block}{Selección con valores lógicos}
\phantomsection\label{selecciuxf3n-con-valores-luxf3gicos}
\begin{itemize}
\item
  También se pueden seleccionar filas (o columnas) usando valores
  lógicos o booleanos,
\item
  Serán seleccionadas las filas (o columnas) que les corresponda un
  valor \texttt{True}, en otro caso no se selecciona, como hemos visto
  con las series de Pandas.
\item
  Esta forma de seleccionar es muy útil, pues permite filtrar los datos
  por el cumplimiento de ciertas condiciones que sean de interés.
\item
  A continuación se seleccionan las columnas que les corresponda
  \texttt{True}
\end{itemize}
\end{block}

\begin{block}{Selección con valores lógicos}
\phantomsection\label{selecciuxf3n-con-valores-luxf3gicos-1}
\begin{Shaded}
\begin{Highlighting}[]
\NormalTok{df.iloc[:,[}\VariableTok{False}\NormalTok{,}\VariableTok{True}\NormalTok{,}\VariableTok{False}\NormalTok{,}\VariableTok{True}\NormalTok{,}\VariableTok{False}\NormalTok{]]}
\end{Highlighting}
\end{Shaded}

\begin{verbatim}
   col2     rcol2
A     9  3.000000
B     0  0.000000
C     2  1.414214
D     1  1.000000
E     7  2.645751
F     3  1.732051
G     4  2.000000
\end{verbatim}
\end{block}

\begin{block}{Filtrado de datos}
\phantomsection\label{filtrado-de-datos}
Una tarea que se presenta a menudo cuando se analizan datos, es
separarlos por niveles o valores de alguna variable, por ejemplo,
género.

\phantomsection\label{annotated-cell-35}%
\begin{Shaded}
\begin{Highlighting}[]
\NormalTok{df.iloc[}\BuiltInTok{list}\NormalTok{(df.col3}\OperatorTok{==}\StringTok{"F"}\NormalTok{)] }\hspace*{\fill}\NormalTok{\circled{1}}
\NormalTok{df.iloc[}\BuiltInTok{list}\NormalTok{(df.col2}\OperatorTok{\textgreater{}=}\DecValTok{7}\NormalTok{)] }\hspace*{\fill}\NormalTok{\circled{2}}
\NormalTok{df.iloc[df.index }\OperatorTok{\%} \DecValTok{2} \OperatorTok{==} \DecValTok{0}\NormalTok{] }\hspace*{\fill}\NormalTok{\circled{3}}
\end{Highlighting}
\end{Shaded}

\begin{description}
\tightlist
\item[\circled{1}]
Seleccionamos del marco de datos \texttt{df} las filas para las cuales
\texttt{col3} corresponde al valor \texttt{F}.
\item[\circled{2}]
Seleccionamos las filas de \texttt{df} que cumplan con la condición que
\texttt{col2} sea mayor o igual que 7.
\item[\circled{3}]
Seleccionemos del marco de datos \texttt{df} las filas que correspondan
a índice par.
\end{description}
\end{block}

\begin{block}{Método \texttt{.loc}}
\phantomsection\label{muxe9todo-.loc}
\begin{itemize}
\item
  Con el método \texttt{loc} se puede acceder a un grupo de filas o
  columnas de la tabla de datos mediante etiquetas, o por medio de un
  arreglo cuyos elementos son valores lógicos.
\item
  Esto permite seleccionar datos basados en los nombres de las etiquetas
  y no en sus posiciones. En el siguiente bloque de código se muestran
  algunos ejemplos
\end{itemize}

\phantomsection\label{annotated-cell-36}%
\begin{Shaded}
\begin{Highlighting}[]
\NormalTok{df.index}\OperatorTok{=}\NormalTok{[}\StringTok{\textquotesingle{}A\textquotesingle{}}\NormalTok{,}\StringTok{\textquotesingle{}B\textquotesingle{}}\NormalTok{,}\StringTok{\textquotesingle{}C\textquotesingle{}}\NormalTok{,}\StringTok{\textquotesingle{}D\textquotesingle{}}\NormalTok{,}\StringTok{\textquotesingle{}E\textquotesingle{}}\NormalTok{,}\StringTok{\textquotesingle{}F\textquotesingle{}}\NormalTok{,}\StringTok{\textquotesingle{}G\textquotesingle{}}\NormalTok{] }\hspace*{\fill}\NormalTok{\circled{1}}
\NormalTok{df.loc[}\StringTok{\textquotesingle{}D\textquotesingle{}}\NormalTok{,}\StringTok{\textquotesingle{}col2\textquotesingle{}}\NormalTok{] }\hspace*{\fill}\NormalTok{\circled{2}}
\NormalTok{df.loc[[}\StringTok{\textquotesingle{}A\textquotesingle{}}\NormalTok{,}\StringTok{\textquotesingle{}D\textquotesingle{}}\NormalTok{]] }\hspace*{\fill}\NormalTok{\circled{3}}
\NormalTok{df.loc[[}\StringTok{\textquotesingle{}A\textquotesingle{}}\NormalTok{,}\StringTok{\textquotesingle{}D\textquotesingle{}}\NormalTok{],[}\StringTok{\textquotesingle{}col1\textquotesingle{}}\NormalTok{,}\StringTok{\textquotesingle{}col3\textquotesingle{}}\NormalTok{]] }\hspace*{\fill}\NormalTok{\circled{4}}
\NormalTok{df.loc[:,[}\StringTok{\textquotesingle{}col1\textquotesingle{}}\NormalTok{,}\StringTok{\textquotesingle{}col3\textquotesingle{}}\NormalTok{]] }\hspace*{\fill}\NormalTok{\circled{5}}
\NormalTok{df.loc[df.col1}\OperatorTok{\textgreater{}}\DecValTok{4}\NormalTok{] }\hspace*{\fill}\NormalTok{\circled{6}}
\NormalTok{df.loc[(df.col1}\OperatorTok{\textgreater{}=}\DecValTok{4}\NormalTok{) }\OperatorTok{\&}\NormalTok{ (df.col2}\OperatorTok{\textless{}=}\DecValTok{3}\NormalTok{)] }\hspace*{\fill}\NormalTok{\circled{7}}
\end{Highlighting}
\end{Shaded}

\begin{description}
\tightlist
\item[\circled{1}]
Asignamos letras como etiquetas de filas.
\item[\circled{2}]
Selección de la fila \texttt{D} y columna \texttt{col2}.
\item[\circled{3}]
Selección de las filas con etiquetas \texttt{A} y \texttt{D}.
\item[\circled{4}]
Selección de las filas con etiquetas \texttt{A} y \texttt{D} y columnas
\texttt{col1} y \texttt{col3}.
\item[\circled{5}]
Selección de las columnas \texttt{col1} y \texttt{col3}, todas las
filas.
\item[\circled{6}]
Selección de aquellas filas que cumplen con la condición que
\texttt{col1} es mayor que 4 y todas las columnas.
\item[\circled{7}]
Selección de las filas que cumplan simultáneamente las condiciones
\texttt{col1\textgreater{}4} y \texttt{col2\textless{}=3}
\end{description}
\end{block}

\begin{block}{Recuerde!}
\phantomsection\label{recuerde}
\begin{tcolorbox}[enhanced jigsaw, opacitybacktitle=0.6, coltitle=black, bottomrule=.15mm, opacityback=0, colback=white, colframe=quarto-callout-note-color-frame, rightrule=.15mm, left=2mm, toptitle=1mm, arc=.35mm, bottomtitle=1mm, titlerule=0mm, title=\textcolor{quarto-callout-note-color}{\faInfo}\hspace{0.5em}{Nota}, toprule=.15mm, colbacktitle=quarto-callout-note-color!10!white, leftrule=.75mm, breakable]

Recuerde que estas operaciones no modifican al data \texttt{df}, si
desea guardar las consultas, deberán asignarse a algún objeto.

\end{tcolorbox}
\end{block}
\end{frame}

\section{Unir Marcos de datos de
Pandas}\label{unir-marcos-de-datos-de-pandas}

\begin{frame}[fragile]{Unir Marcos de datos de Pandas}
\begin{block}{Unir marcos de datos de Pandas}
\phantomsection\label{unir-marcos-de-datos-de-pandas-1}
\index{Pandas DataFrame!Función!concat}

En caso que se tengan dos marcos de datos con exactamente las mismas
columnas, por ejemplo, cuando tenemos información de hombres y mujeres
separada en objetos distintos, podríamos unir los dos en uno, con la
función \index{Función!concat} \texttt{concat}.
\end{block}

\begin{block}{Marcos de datos para unir}
\phantomsection\label{marcos-de-datos-para-unir}
\begin{itemize}
\item
  Para ilustrar el uso de \texttt{concat} crearemos dos marcos de datos,
  \texttt{dfM} y \texttt{dfF} que contienen datos de edad, estatura, y
  género.
\item
  Note que todas las filas de \texttt{dfM} tienen \texttt{Genero=M} y
  todas las de \texttt{dfF}, tienen \texttt{Genero=F}.
\end{itemize}

\begin{Shaded}
\begin{Highlighting}[]
\NormalTok{dfM}\OperatorTok{=}\NormalTok{pd.DataFrame(\{}\StringTok{"Edad"}\NormalTok{:[}\DecValTok{24}\NormalTok{,}\DecValTok{21}\NormalTok{,}\DecValTok{18}\NormalTok{,}\DecValTok{26}\NormalTok{],}
     \StringTok{"Estat"}\NormalTok{:[}\FloatTok{1.65}\NormalTok{,}\FloatTok{1.72}\NormalTok{,}\FloatTok{1.69}\NormalTok{,}\FloatTok{1.8}\NormalTok{],}
     \StringTok{"Genero"}\NormalTok{:[}\StringTok{"M"}\NormalTok{,}\StringTok{"M"}\NormalTok{,}\StringTok{"M"}\NormalTok{,}\StringTok{"M"}\NormalTok{]\})}
\NormalTok{dfF}\OperatorTok{=}\NormalTok{pd.DataFrame(\{}\StringTok{"Estat"}\NormalTok{:[}\FloatTok{1.60}\NormalTok{,}\FloatTok{1.58}\NormalTok{,}\FloatTok{1.62}\NormalTok{,}\FloatTok{1.71}\NormalTok{,}\FloatTok{1.64}\NormalTok{],}
\StringTok{"Genero"}\NormalTok{:[}\StringTok{"F"}\NormalTok{,}\StringTok{"F"}\NormalTok{,}\StringTok{"F"}\NormalTok{,}\StringTok{"F"}\NormalTok{,}\StringTok{"F"}\NormalTok{],     }
\StringTok{"Edad"}\NormalTok{:[}\DecValTok{23}\NormalTok{,}\DecValTok{22}\NormalTok{,}\DecValTok{19}\NormalTok{,}\DecValTok{25}\NormalTok{,}\DecValTok{21}\NormalTok{]\})}
\end{Highlighting}
\end{Shaded}
\end{block}

\begin{block}{Uso de \texttt{concat}}
\phantomsection\label{uso-de-concat}
Usemos la función \texttt{concat()} para unir los dos marcos de datos en
uno solo, lo haremos uno debajo del otro

\phantomsection\label{annotated-cell-38}%
\begin{Shaded}
\begin{Highlighting}[]
\NormalTok{tablas}\OperatorTok{=}\NormalTok{[dfM,dfF] }\hspace*{\fill}\NormalTok{\circled{1}}
\NormalTok{df1}\OperatorTok{=}\NormalTok{pd.concat(tablas) }\hspace*{\fill}\NormalTok{\circled{2}}
\NormalTok{df1}
\end{Highlighting}
\end{Shaded}

\begin{description}
\tightlist
\item[\circled{1}]
Lista con los \texttt{DataFrames} de pandas.
\item[\circled{2}]
Unimos, un marco de datos debajo del otro.
\end{description}

\begin{verbatim}
   Edad  Estat Genero
0    24   1.65      M
1    21   1.72      M
2    18   1.69      M
3    26   1.80      M
0    23   1.60      F
1    22   1.58      F
2    19   1.62      F
3    25   1.71      F
4    21   1.64      F
\end{verbatim}
\end{block}

\begin{block}{Reestablecer índices}
\phantomsection\label{reestablecer-uxedndices}
Notemos, en la salida anterior, que los índices de filas son los mismos
de los marcos de datos originales, es decir, tenemos índices repetidos,
esto puede causar confusión al momento de seleccionar filas usando
índices, para evitar ese comportamiento usamos la opción
\texttt{ignore\_index=True}

\begin{Shaded}
\begin{Highlighting}[]
\NormalTok{df1}\OperatorTok{=}\NormalTok{pd.concat(tablas,ignore\_index}\OperatorTok{=}\VariableTok{True}\NormalTok{)}
\NormalTok{df1}
\end{Highlighting}
\end{Shaded}

\begin{verbatim}
   Edad  Estat Genero
0    24   1.65      M
1    21   1.72      M
2    18   1.69      M
3    26   1.80      M
4    23   1.60      F
5    22   1.58      F
6    19   1.62      F
7    25   1.71      F
8    21   1.64      F
\end{verbatim}
\end{block}

\begin{block}{Claves para los datas originales}
\phantomsection\label{claves-para-los-datas-originales}
Es posible asignar claves a cada bloque representado, en este caso, por
los niveles de género, para que en caso que sea necesario, se facilite
la posterior separación

\begin{Shaded}
\begin{Highlighting}[]
\NormalTok{dfC}\OperatorTok{=}\NormalTok{pd.concat(tablas,keys}\OperatorTok{=}\NormalTok{[}\StringTok{\textquotesingle{}Masculino\textquotesingle{}}\NormalTok{, }\StringTok{\textquotesingle{}Femenino\textquotesingle{}}\NormalTok{])}
\NormalTok{dfC}
\end{Highlighting}
\end{Shaded}

\begin{verbatim}
             Edad  Estat Genero
Masculino 0    24   1.65      M
          1    21   1.72      M
          2    18   1.69      M
          3    26   1.80      M
Femenino  0    23   1.60      F
          1    22   1.58      F
          2    19   1.62      F
          3    25   1.71      F
          4    21   1.64      F
\end{verbatim}
\end{block}

\begin{block}{Separación usando las claves}
\phantomsection\label{separaciuxf3n-usando-las-claves}
De ese modo, si queremos recuperar el grupo femenino (\texttt{F})
podríamos hacerlo de la siguiente forma.

\begin{Shaded}
\begin{Highlighting}[]
\NormalTok{dfC.loc[}\StringTok{\textquotesingle{}Femenino\textquotesingle{}}\NormalTok{]}
\end{Highlighting}
\end{Shaded}

\begin{verbatim}
   Edad  Estat Genero
0    23   1.60      F
1    22   1.58      F
2    19   1.62      F
3    25   1.71      F
4    21   1.64      F
\end{verbatim}
\end{block}

\begin{block}{Multi index}
\phantomsection\label{multi-index}
\begin{Shaded}
\begin{Highlighting}[]
\NormalTok{df1.index}
\end{Highlighting}
\end{Shaded}

\begin{verbatim}
RangeIndex(start=0, stop=9, step=1)
\end{verbatim}

\begin{Shaded}
\begin{Highlighting}[]
\NormalTok{dfC.index}
\end{Highlighting}
\end{Shaded}

\begin{verbatim}
MultiIndex([('Masculino', 0),
            ('Masculino', 1),
            ('Masculino', 2),
            ('Masculino', 3),
            ( 'Femenino', 0),
            ( 'Femenino', 1),
            ( 'Femenino', 2),
            ( 'Femenino', 3),
            ( 'Femenino', 4)],
           )
\end{verbatim}
\end{block}

\begin{block}{El argumento \texttt{join}}
\phantomsection\label{el-argumento-join}
Cuando los marcos de datos que queremos concatenar tienen algunas
columnas que no son comunes, \texttt{concat()} hace coincidir las
columnas comunes y las no comunes las adiciona a la tabla y llena las
filas que quedan sin datos con \texttt{NaN}, veamos un ejemplo
\end{block}

\begin{block}{Datos para ilustrar el argumento \texttt{join}}
\phantomsection\label{datos-para-ilustrar-el-argumento-join}
\begin{Shaded}
\begin{Highlighting}[]
\NormalTok{df1 }\OperatorTok{=}\NormalTok{ pd.DataFrame(\{}\StringTok{\textquotesingle{}A\textquotesingle{}}\NormalTok{: [}\StringTok{\textquotesingle{}A0\textquotesingle{}}\NormalTok{, }\StringTok{\textquotesingle{}A1\textquotesingle{}}\NormalTok{, }\StringTok{\textquotesingle{}A2\textquotesingle{}}\NormalTok{, }\StringTok{\textquotesingle{}A3\textquotesingle{}}\NormalTok{],}
                    \StringTok{\textquotesingle{}B\textquotesingle{}}\NormalTok{: [}\StringTok{\textquotesingle{}B0\textquotesingle{}}\NormalTok{, }\StringTok{\textquotesingle{}B1\textquotesingle{}}\NormalTok{, }\StringTok{\textquotesingle{}B2\textquotesingle{}}\NormalTok{, }\StringTok{\textquotesingle{}B3\textquotesingle{}}\NormalTok{],}
                    \StringTok{\textquotesingle{}C\textquotesingle{}}\NormalTok{: [}\StringTok{\textquotesingle{}C0\textquotesingle{}}\NormalTok{, }\StringTok{\textquotesingle{}C1\textquotesingle{}}\NormalTok{, }\StringTok{\textquotesingle{}C2\textquotesingle{}}\NormalTok{, }\StringTok{\textquotesingle{}C3\textquotesingle{}}\NormalTok{],}
                    \StringTok{\textquotesingle{}D\textquotesingle{}}\NormalTok{: [}\StringTok{\textquotesingle{}D0\textquotesingle{}}\NormalTok{, }\StringTok{\textquotesingle{}D1\textquotesingle{}}\NormalTok{, }\StringTok{\textquotesingle{}D2\textquotesingle{}}\NormalTok{, }\StringTok{\textquotesingle{}D3\textquotesingle{}}\NormalTok{]\},}
\NormalTok{                       index}\OperatorTok{=}\NormalTok{[}\DecValTok{0}\NormalTok{, }\DecValTok{1}\NormalTok{, }\DecValTok{2}\NormalTok{, }\DecValTok{3}\NormalTok{])}
\NormalTok{df1}
\end{Highlighting}
\end{Shaded}

\begin{verbatim}
    A   B   C   D
0  A0  B0  C0  D0
1  A1  B1  C1  D1
2  A2  B2  C2  D2
3  A3  B3  C3  D3
\end{verbatim}

Observe que \texttt{df1} tiene cuatro columnas, nombradas \texttt{A},
\texttt{B}, \texttt{C} y \texttt{D}.
\end{block}

\begin{block}{Datos para ilustrar el argumento \texttt{join} (II)}
\phantomsection\label{datos-para-ilustrar-el-argumento-join-ii}
Ahora creamos el marco de datos \texttt{df2} que tiene tres columnas, de
las cuales dos tienen nombres que coinciden con las de \texttt{df1}, es
decir, \texttt{B} y \texttt{D}, son nombres de columnas comunes a las
del marco de datos \texttt{df1}.

\begin{Shaded}
\begin{Highlighting}[]
\NormalTok{df2 }\OperatorTok{=}\NormalTok{ pd.DataFrame(\{}\StringTok{\textquotesingle{}B\textquotesingle{}}\NormalTok{: [}\StringTok{\textquotesingle{}B2\textquotesingle{}}\NormalTok{, }\StringTok{\textquotesingle{}B3\textquotesingle{}}\NormalTok{, }\StringTok{\textquotesingle{}B6\textquotesingle{}}\NormalTok{, }\StringTok{\textquotesingle{}B7\textquotesingle{}}\NormalTok{],}
                    \StringTok{\textquotesingle{}D\textquotesingle{}}\NormalTok{: [}\StringTok{\textquotesingle{}D2\textquotesingle{}}\NormalTok{, }\StringTok{\textquotesingle{}D3\textquotesingle{}}\NormalTok{, }\StringTok{\textquotesingle{}D6\textquotesingle{}}\NormalTok{, }\StringTok{\textquotesingle{}D7\textquotesingle{}}\NormalTok{],}
                    \StringTok{\textquotesingle{}F\textquotesingle{}}\NormalTok{: [}\StringTok{\textquotesingle{}F2\textquotesingle{}}\NormalTok{, }\StringTok{\textquotesingle{}F3\textquotesingle{}}\NormalTok{, }\StringTok{\textquotesingle{}F6\textquotesingle{}}\NormalTok{, }\StringTok{\textquotesingle{}F7\textquotesingle{}}\NormalTok{]\},}
\NormalTok{                         index}\OperatorTok{=}\NormalTok{[}\DecValTok{2}\NormalTok{, }\DecValTok{3}\NormalTok{, }\DecValTok{6}\NormalTok{, }\DecValTok{7}\NormalTok{])}
\NormalTok{df2}
\end{Highlighting}
\end{Shaded}

\begin{verbatim}
    B   D   F
2  B2  D2  F2
3  B3  D3  F3
6  B6  D6  F6
7  B7  D7  F7
\end{verbatim}
\end{block}

\begin{block}{Datos para ilustrar el argumento \texttt{join} (III)}
\phantomsection\label{datos-para-ilustrar-el-argumento-join-iii}
Fíjese, además, que \texttt{df1} y \texttt{df2} tienen dos filas en
común, cuyos índices son \texttt{2} y \texttt{3}. Ahora concatenamos

\begin{Shaded}
\begin{Highlighting}[]
\NormalTok{dfC12}\OperatorTok{=}\NormalTok{pd.concat([df1,df2])   }
\NormalTok{dfC12}
\end{Highlighting}
\end{Shaded}

\begin{verbatim}
     A   B    C   D    F
0   A0  B0   C0  D0  NaN
1   A1  B1   C1  D1  NaN
2   A2  B2   C2  D2  NaN
3   A3  B3   C3  D3  NaN
2  NaN  B2  NaN  D2   F2
3  NaN  B3  NaN  D3   F3
6  NaN  B6  NaN  D6   F6
7  NaN  B7  NaN  D7   F7
\end{verbatim}

Observe que se hacen coincidir las columnas cuidando que la información
de las filas con el mismo indice se mantenga.
\end{block}

\begin{block}{Uso de \texttt{join="inner"}}
\phantomsection\label{uso-de-joininner}
Con la opción \texttt{join="inner"} solo se concatenan las columnas que
sean comunes a los dos marcos de datos

\begin{Shaded}
\begin{Highlighting}[]
\NormalTok{dfC12}\OperatorTok{=}\NormalTok{pd.concat([df1,df2],join}\OperatorTok{=}\StringTok{"inner"}\NormalTok{)}
\NormalTok{dfC12}
\end{Highlighting}
\end{Shaded}

\begin{verbatim}
    B   D
0  B0  D0
1  B1  D1
2  B2  D2
3  B3  D3
2  B2  D2
3  B3  D3
6  B6  D6
7  B7  D7
\end{verbatim}

La opción \texttt{axis=1} es para concatenar por filas, es decir coloca
un marco de datos al lado del otro, haciendo coincidir las filas con el
mismo índice

\begin{Shaded}
\begin{Highlighting}[]
\NormalTok{dfC12}\OperatorTok{=}\NormalTok{pd.concat([df1,df2],axis}\OperatorTok{=}\DecValTok{1}\NormalTok{) }
\NormalTok{dfC12}
\end{Highlighting}
\end{Shaded}

\begin{verbatim}
     A    B    C    D    B    D    F
0   A0   B0   C0   D0  NaN  NaN  NaN
1   A1   B1   C1   D1  NaN  NaN  NaN
2   A2   B2   C2   D2   B2   D2   F2
3   A3   B3   C3   D3   B3   D3   F3
6  NaN  NaN  NaN  NaN   B6   D6   F6
7  NaN  NaN  NaN  NaN   B7   D7   F7
\end{verbatim}
\end{block}

\begin{block}{Concatenar datos por filas}
\phantomsection\label{concatenar-datos-por-filas}
De manera análoga que con las columnas, con la opción
\texttt{join="inner"} solo se conservan las filas que sean comunes a los
dos marcos de datos

\begin{Shaded}
\begin{Highlighting}[]
\NormalTok{dfC12}\OperatorTok{=}\NormalTok{pd.concat([df1,df2],axis}\OperatorTok{=}\DecValTok{1}\NormalTok{,join}\OperatorTok{=}\StringTok{"inner"}\NormalTok{)}
\NormalTok{dfC12}
\end{Highlighting}
\end{Shaded}

\begin{verbatim}
    A   B   C   D   B   D   F
2  A2  B2  C2  D2  B2  D2  F2
3  A3  B3  C3  D3  B3  D3  F3
\end{verbatim}
\end{block}

\begin{block}{Ordenar un marco de datos}
\phantomsection\label{ordenar-un-marco-de-datos}
Para ordenar un \texttt{DataFrame} de Pandas se cuenta con el método
\index{Pandas DataFrame!Método!sort\_values} \texttt{sort\_values}. En
los siguientes ejemplos se usa el marco de datos \texttt{df}, el cual
ordenamos por la columna \texttt{col1}

\begin{Shaded}
\begin{Highlighting}[]
\NormalTok{df\_or}\OperatorTok{=}\NormalTok{df.sort\_values(by}\OperatorTok{=}\NormalTok{[}\StringTok{\textquotesingle{}col1\textquotesingle{}}\NormalTok{]) }\CommentTok{\#  }
\NormalTok{df\_or }\CommentTok{\# df ordenado  }
\end{Highlighting}
\end{Shaded}

\begin{verbatim}
   col1  col2 col3     rcol2  sum_col
A     1     9    F  3.000000       10
C     2     2    F  1.414214        4
B     3     0    M  0.000000        3
D     4     1    F  1.000000        5
G     5     4    M  2.000000        9
E     7     7    F  2.645751       14
F     9     3    M  1.732051       12
\end{verbatim}
\end{block}

\begin{block}{Orenar por más de una columna}
\phantomsection\label{orenar-por-muxe1s-de-una-columna}
También es posible ordenar por más de una columna, por ejemplo,
ordenamos \texttt{df} primero por \texttt{col3} y luego por
\texttt{col1}

\begin{Shaded}
\begin{Highlighting}[]
\NormalTok{df\_or}\OperatorTok{=}\NormalTok{df.sort\_values(by}\OperatorTok{=}\NormalTok{[}\StringTok{\textquotesingle{}col3\textquotesingle{}}\NormalTok{,}\StringTok{\textquotesingle{}col1\textquotesingle{}}\NormalTok{])}
\NormalTok{df\_or }\CommentTok{\# df ordenado }
\end{Highlighting}
\end{Shaded}

\begin{verbatim}
   col1  col2 col3     rcol2  sum_col
A     1     9    F  3.000000       10
C     2     2    F  1.414214        4
D     4     1    F  1.000000        5
E     7     7    F  2.645751       14
B     3     0    M  0.000000        3
G     5     4    M  2.000000        9
F     9     3    M  1.732051       12
\end{verbatim}
\end{block}

\begin{block}{Variables categóricas con pandas}
\phantomsection\label{sec-VariablesCatPandas}
\index{Variables!Categóricas}\index{Pandas DataFrame!Variables!Categóricas}

\begin{itemize}
\item
  En el análisis de datos, con frecuencia, es necesario tratar con
  variables categóricas o factores.
\item
  Estudiaremos como hacer la conversión de columnas de un \textbf{data
  frame} de pandas, cuyo contenido es texto \texttt{(object)} o
  numérico, a categórico (\texttt{category}).
\item
  Para los ejemplos usaremos el \textbf{data frame} \texttt{df\_es} que
  hemos leído
\item
  A continuación leemos los datos que se encuentran en un archivo
  \texttt{csv}
\end{itemize}
\end{block}

\begin{block}{Lectura de los datos}
\phantomsection\label{lectura-de-los-datos}
\begin{Shaded}
\begin{Highlighting}[]
\NormalTok{df\_es}\OperatorTok{=}\NormalTok{pd.read\_csv(}\StringTok{\textquotesingle{}DatosEst.csv\textquotesingle{}}\NormalTok{)}
\end{Highlighting}
\end{Shaded}

Una forma de conocer el tipo de datos que contienen las columnas
\texttt{sexo}, \texttt{Tcolegio} y \texttt{Afisica}, es mediante el
método \texttt{.info}, como se ilustra a continuación

\begin{Shaded}
\begin{Highlighting}[]
\NormalTok{df\_es[[}\StringTok{\textquotesingle{}Sexo\textquotesingle{}}\NormalTok{,}\StringTok{\textquotesingle{}Tcolegio\textquotesingle{}}\NormalTok{,}\StringTok{\textquotesingle{}Afisica\textquotesingle{}}\NormalTok{]].info()}
\end{Highlighting}
\end{Shaded}

\begin{verbatim}
<class 'pandas.core.frame.DataFrame'>
RangeIndex: 45 entries, 0 to 44
Data columns (total 3 columns):
 #   Column    Non-Null Count  Dtype 
---  ------    --------------  ----- 
 0   Sexo      45 non-null     object
 1   Tcolegio  45 non-null     object
 2   Afisica   45 non-null     object
dtypes: object(3)
memory usage: 1.2+ KB
\end{verbatim}
\end{block}

\begin{block}{Conversión con \texttt{.astype("category")}}
\phantomsection\label{conversiuxf3n-con-.astypecategory}
\begin{itemize}
\item
  La columnas, que guardan información de variables categóricas,
  aparecen como de tipo \texttt{object}
\item
  Esta es la forma como pandas identifica el contenido de texto o
  \texttt{string}.
\item
  A continuación convertimos dichas columnas a tipo \texttt{category}
\end{itemize}

\begin{Shaded}
\begin{Highlighting}[]
\NormalTok{df\_es[}\StringTok{\textquotesingle{}Sexo\textquotesingle{}}\NormalTok{]}\OperatorTok{=}\NormalTok{df\_es[}\StringTok{\textquotesingle{}Sexo\textquotesingle{}}\NormalTok{].astype(}\StringTok{"category"}\NormalTok{)}
\NormalTok{df\_es[}\StringTok{\textquotesingle{}Tcolegio\textquotesingle{}}\NormalTok{]}\OperatorTok{=}\NormalTok{df\_es[}\StringTok{\textquotesingle{}Tcolegio\textquotesingle{}}\NormalTok{].astype(}\StringTok{"category"}\NormalTok{)}
\NormalTok{df\_es[}\StringTok{\textquotesingle{}Afisica\textquotesingle{}}\NormalTok{]}\OperatorTok{=}\NormalTok{df\_es[}\StringTok{\textquotesingle{}Afisica\textquotesingle{}}\NormalTok{].astype(}\StringTok{"category"}\NormalTok{)}
\end{Highlighting}
\end{Shaded}
\end{block}

\begin{block}{Tabla de contingencia}
\phantomsection\label{tabla-de-contingencia}
\begin{Shaded}
\begin{Highlighting}[]
\NormalTok{pd.crosstab(df\_es.Sexo,df\_es.Tcolegio) }
\end{Highlighting}
\end{Shaded}

\begin{verbatim}
Tcolegio  Privado  Público
Sexo                      
Hombre          6       16
Mujer           3       20
\end{verbatim}
\end{block}

\begin{block}{Información y renombrado de categorías}
\phantomsection\label{informaciuxf3n-y-renombrado-de-categoruxedas}
Ahora las columnas \texttt{Sexo}, \texttt{Tcolegio} y \texttt{Afisica}
son columnas categóricas, es decir, de tipo \texttt{category}. Podemos
pedir que se imprima la información de las categorías de una variable
categórica, invocando el método \texttt{cat.categories}, y renombrar las
categorías con \texttt{cat.rename\_categories}. A continuación lo
hacemos con la columna \texttt{Sexo}
\end{block}

\begin{block}{Información y renombrado de categorías (II)}
\phantomsection\label{informaciuxf3n-y-renombrado-de-categoruxedas-ii}
\phantomsection\label{annotated-cell-56}%
\begin{Shaded}
\begin{Highlighting}[]
\NormalTok{df\_es[}\StringTok{\textquotesingle{}Sexo\textquotesingle{}}\NormalTok{].cat.categories }\hspace*{\fill}\NormalTok{\circled{1}}
\NormalTok{n\_cat}\OperatorTok{=}\NormalTok{[}\StringTok{\textquotesingle{}Grupo Hombres\textquotesingle{}}\NormalTok{, }\StringTok{\textquotesingle{}Grupo Mujeres\textquotesingle{}}\NormalTok{] }\hspace*{\fill}\NormalTok{\circled{2}}
\NormalTok{df\_es[}\StringTok{\textquotesingle{}Sexo\textquotesingle{}}\NormalTok{]}\OperatorTok{=}\NormalTok{df\_es[}\StringTok{\textquotesingle{}Sexo\textquotesingle{}}\NormalTok{].cat.rename\_categories(n\_cat) }\hspace*{\fill}\NormalTok{\circled{3}}
\NormalTok{df\_es[}\StringTok{\textquotesingle{}Sexo\textquotesingle{}}\NormalTok{].cat.categories}
\end{Highlighting}
\end{Shaded}

\begin{description}
\tightlist
\item[\circled{1}]
Se recuperan las categorías de la columna categórica \texttt{Sexo}.
\item[\circled{2}]
Para renombrar las categorías, creamos una lista que contiene los nuevos
nombres de las categorías.
\item[\circled{3}]
Renombramos las categorías.
\end{description}

\begin{verbatim}
Index(['Hombre', 'Mujer'], dtype='object')
Index(['Grupo Hombres', 'Grupo Mujeres'], dtype='object')
\end{verbatim}
\end{block}

\begin{block}{Información y renombrado de categorías (II)}
\phantomsection\label{informaciuxf3n-y-renombrado-de-categoruxedas-ii-1}
\phantomsection\label{annotated-cell-57}%
\begin{Shaded}
\begin{Highlighting}[]
\NormalTok{df\_es[}\StringTok{\textquotesingle{}Sexo\textquotesingle{}}\NormalTok{].cat.categories }\hspace*{\fill}\NormalTok{\circled{1}}
\NormalTok{n\_cat}\OperatorTok{=}\NormalTok{[}\StringTok{\textquotesingle{}Grupo Hombres\textquotesingle{}}\NormalTok{, }\StringTok{\textquotesingle{}Grupo Mujeres\textquotesingle{}}\NormalTok{] }\hspace*{\fill}\NormalTok{\circled{2}}
\NormalTok{df\_es[}\StringTok{\textquotesingle{}Sexo\textquotesingle{}}\NormalTok{]}\OperatorTok{=}\NormalTok{df\_es[}\StringTok{\textquotesingle{}Sexo\textquotesingle{}}\NormalTok{].cat.rename\_categories(n\_cat) }\hspace*{\fill}\NormalTok{\circled{3}}
\end{Highlighting}
\end{Shaded}

\begin{verbatim}
Index(['Grupo Hombres', 'Grupo Mujeres'], dtype='object')
\end{verbatim}

escribe la columna \texttt{Sexo}, si se desea conservar la columna
original, debemos crear la columna con otro nombre, algo como
\texttt{df\_es{[}\textquotesingle{}SexoR\textquotesingle{}{]}=...} 1.
Otra forma de hacer exactamente lo mismo es mediante un diccionario que
contiene, como claves, las categorías actuales y como valores los nuevos
nombres. 2. Se renombran las categorías, como en 3. 3. Se imprime como
quedaron los nombres de las categorías.
\end{block}

\begin{block}{Información de actividad física}
\phantomsection\label{informaciuxf3n-de-actividad-fuxedsica}
\begin{Shaded}
\begin{Highlighting}[]
\NormalTok{df\_es[}\StringTok{\textquotesingle{}Afisica\textquotesingle{}}\NormalTok{].cat.categories }
\end{Highlighting}
\end{Shaded}

\begin{verbatim}
Index(['Cinco días a la semana', 'Cuatro días a la semana',
       'Dos días a la semana', 'Ningún día a la semana',
       'Todos los días de la semana', 'Tres días a la semana'],
      dtype='object')
\end{verbatim}
\end{block}

\begin{block}{Recategorización}
\phantomsection\label{recategorizaciuxf3n}
Una tarea frecuente con variables categóricas consiste en hacer una
recategorización de sus niveles, por ejemplo, suponga que a partir de la
columna \texttt{Afisica} queremos construir una nueva variable
categórica, con dos categorías, de tal forma que los niveles
\texttt{Cuatro\ días\ a\ la\ semana},
\texttt{Cinco\ días\ a\ la\ semana} y
\texttt{Todos\ los\ días\ de\ la\ semana} conformen la categoría
\texttt{Alta}; \texttt{Ningún\ día\ a\ la\ semana},
\texttt{Dos\ días\ a\ la\ semana} y \texttt{Tres\ días\ a\ la\ semana},
conformen la categoría \texttt{Baja}.
\end{block}

\begin{block}{Recategorización (II)}
\phantomsection\label{recategorizaciuxf3n-ii}
Lo anterior se puede llevar a cabo mediante un \texttt{mapeo}, de la
siguiente forma:

\phantomsection\label{annotated-cell-59}%
\begin{Shaded}
\begin{Highlighting}[]
\NormalTok{dic\_map}\OperatorTok{=}\NormalTok{\{}\StringTok{\textquotesingle{}Ningún día a la semana\textquotesingle{}}\NormalTok{:}\StringTok{\textquotesingle{}Baja\textquotesingle{}}\NormalTok{,}\StringTok{\textquotesingle{}Dos días a la semana\textquotesingle{}}\NormalTok{:}\StringTok{\textquotesingle{}Baja\textquotesingle{}}\NormalTok{,}
\StringTok{\textquotesingle{}Tres días a la semana\textquotesingle{}}\NormalTok{:}\StringTok{\textquotesingle{}Baja\textquotesingle{}}\NormalTok{,}
\StringTok{\textquotesingle{}Todos los días de la semana\textquotesingle{}}\NormalTok{:}\StringTok{\textquotesingle{}Alta\textquotesingle{}}\NormalTok{,}
\StringTok{\textquotesingle{}Cinco días a la semana\textquotesingle{}}\NormalTok{:}\StringTok{\textquotesingle{}Alta\textquotesingle{}}\NormalTok{,}
\StringTok{\textquotesingle{}Cuatro días a la semana\textquotesingle{}}\NormalTok{:}\StringTok{\textquotesingle{}Alta\textquotesingle{}}\NormalTok{\} }\hspace*{\fill}\NormalTok{\circled{1}}
\NormalTok{df\_es[}\StringTok{\textquotesingle{}Afisica2C\textquotesingle{}}\NormalTok{]}\OperatorTok{=}\NormalTok{df\_es[}\StringTok{\textquotesingle{}Afisica\textquotesingle{}}\NormalTok{].}\BuiltInTok{map}\NormalTok{(dic\_map).astype(}\StringTok{"category"}\NormalTok{) }\hspace*{\fill}\NormalTok{\circled{2}}

\NormalTok{df\_es[}\StringTok{\textquotesingle{}Afisica2C\textquotesingle{}}\NormalTok{].value\_counts()}
\end{Highlighting}
\end{Shaded}

\begin{verbatim}
Afisica2C
Baja    40
Alta     5
Name: count, dtype: int64
\end{verbatim}
\end{block}

\begin{block}{Recategorización (III)}
\phantomsection\label{recategorizaciuxf3n-iii}
\phantomsection\label{annotated-cell-60}%
\begin{Shaded}
\begin{Highlighting}[]
\NormalTok{dic\_map}\OperatorTok{=}\NormalTok{\{}\StringTok{\textquotesingle{}Ningún día a la semana\textquotesingle{}}\NormalTok{:}\StringTok{\textquotesingle{}Baja\textquotesingle{}}\NormalTok{,}
  \StringTok{\textquotesingle{}Dos días a la semana\textquotesingle{}}\NormalTok{:}\StringTok{\textquotesingle{}Baja\textquotesingle{}}\NormalTok{,}
  \StringTok{\textquotesingle{}Tres días a la semana\textquotesingle{}}\NormalTok{:}\StringTok{\textquotesingle{}Baja\textquotesingle{}}\NormalTok{,}
  \StringTok{\textquotesingle{}Todos los días de la semana\textquotesingle{}}\NormalTok{:}\StringTok{\textquotesingle{}Alta\textquotesingle{}}\NormalTok{,}
  \StringTok{\textquotesingle{}Cinco días a la semana\textquotesingle{}}\NormalTok{:}\StringTok{\textquotesingle{}Alta\textquotesingle{}}\NormalTok{,}
  \StringTok{\textquotesingle{}Cuatro días a la semana\textquotesingle{}}\NormalTok{:}\StringTok{\textquotesingle{}Alta\textquotesingle{}}\NormalTok{\} }\hspace*{\fill}\NormalTok{\circled{1}}
\NormalTok{df\_es[}\StringTok{\textquotesingle{}Afisica2C\textquotesingle{}}\NormalTok{]}\OperatorTok{=}\NormalTok{df\_es[}\StringTok{\textquotesingle{}Afisica\textquotesingle{}}\NormalTok{].}\BuiltInTok{map}\NormalTok{(dic\_map) }\hspace*{\fill}\NormalTok{\circled{2}}
\NormalTok{df\_es[}\StringTok{\textquotesingle{}Afisica2C\textquotesingle{}}\NormalTok{].astype(}\StringTok{"category"}\NormalTok{) }\hspace*{\fill}\NormalTok{\circled{3}}
\end{Highlighting}
\end{Shaded}

\begin{description}
\tightlist
\item[\circled{1}]
Diccionario con los nombres de las categorías actuales como claves y los
nombres de las nuevas categorías como valores.
\item[\circled{2}]
Reemplazamos las antiguas categorías por las nuevas, creando una nueva
columna llamada \texttt{Afisica2C}.
\item[\circled{3}]
Convertimos la nueva columna a categórica.
\end{description}
\end{block}

\begin{block}{Columnas numéricas a factores}
\phantomsection\label{columnas-numuxe9ricas-a-factores}
Cuando organizamos datos para hacer resúmenes o análisis estadístico,
una tarea frecuente es transformar una columna, que inicialmente se
registra como numérica, a una variable categórica o \texttt{factor}.
Consideremos el siguiente marco de datos de pandas, que tiene una
columna numérica y una de texto

\begin{Shaded}
\begin{Highlighting}[]
\NormalTok{d}\OperatorTok{=}\NormalTok{\{ }\StringTok{\textquotesingle{}A\textquotesingle{}}\NormalTok{:[}\DecValTok{1}\NormalTok{,}\DecValTok{4}\NormalTok{,}\DecValTok{2}\NormalTok{,}\DecValTok{1}\NormalTok{,}\DecValTok{3}\NormalTok{,}\DecValTok{2}\NormalTok{,}\DecValTok{4}\NormalTok{,}\DecValTok{3}\NormalTok{,}\DecValTok{2}\NormalTok{],}
    \StringTok{\textquotesingle{}B\textquotesingle{}}\NormalTok{:[}\StringTok{\textquotesingle{}M\textquotesingle{}}\NormalTok{,}\StringTok{\textquotesingle{}F\textquotesingle{}}\NormalTok{,}\StringTok{\textquotesingle{}M\textquotesingle{}}\NormalTok{,}\StringTok{\textquotesingle{}F\textquotesingle{}}\NormalTok{,}\StringTok{\textquotesingle{}M\textquotesingle{}}\NormalTok{,}\StringTok{\textquotesingle{}F\textquotesingle{}}\NormalTok{,}\StringTok{\textquotesingle{}M\textquotesingle{}}\NormalTok{,}\StringTok{\textquotesingle{}F\textquotesingle{}}\NormalTok{,}\StringTok{\textquotesingle{}M\textquotesingle{}}\NormalTok{]\}}
\NormalTok{dfc}\OperatorTok{=}\NormalTok{pd.DataFrame(d) }
\NormalTok{dfc.info()}
\end{Highlighting}
\end{Shaded}

\begin{verbatim}
<class 'pandas.core.frame.DataFrame'>
RangeIndex: 9 entries, 0 to 8
Data columns (total 2 columns):
 #   Column  Non-Null Count  Dtype 
---  ------  --------------  ----- 
 0   A       9 non-null      int64 
 1   B       9 non-null      object
dtypes: int64(1), object(1)
memory usage: 272.0+ bytes
\end{verbatim}
\end{block}

\begin{block}{Columnas numéricas a factores (II)}
\phantomsection\label{columnas-numuxe9ricas-a-factores-ii}
Queremos convertir la columna \texttt{A}, que es de tipo \texttt{int64},
a categórica; de tal forma que los valores 1 pasen a ser \texttt{c1},
los 2 se conviertan en \texttt{c2} y así, hasta el valor 4.

\phantomsection\label{annotated-cell-62}%
\begin{Shaded}
\begin{Highlighting}[]
\NormalTok{rep}\OperatorTok{=}\NormalTok{\{}\DecValTok{1}\NormalTok{:}\StringTok{"c1"}\NormalTok{,}\DecValTok{2}\NormalTok{:}\StringTok{"c2"}\NormalTok{,}\DecValTok{3}\NormalTok{:}\StringTok{"c3"}\NormalTok{,}\DecValTok{4}\NormalTok{:}\StringTok{"c4"}\NormalTok{\} }\hspace*{\fill}\NormalTok{\circled{1}}
\CommentTok{\#dfc[\textquotesingle{}Acat\textquotesingle{}]=dfc[\textquotesingle{}A\textquotesingle{}].replace(rep) }\hspace*{\fill}\NormalTok{\circled{2}}
\NormalTok{dfc[}\StringTok{\textquotesingle{}Acat\textquotesingle{}}\NormalTok{]}\OperatorTok{=}\NormalTok{dfc[}\StringTok{\textquotesingle{}A\textquotesingle{}}\NormalTok{].}\BuiltInTok{map}\NormalTok{(rep) }
\NormalTok{dfc[}\StringTok{\textquotesingle{}Acat\textquotesingle{}}\NormalTok{]}\OperatorTok{=}\NormalTok{dfc[}\StringTok{\textquotesingle{}Acat\textquotesingle{}}\NormalTok{].astype(}\StringTok{\textquotesingle{}category\textquotesingle{}}\NormalTok{) }\hspace*{\fill}\NormalTok{\circled{3}}
\NormalTok{dfc.info() }
\end{Highlighting}
\end{Shaded}

\begin{description}
\tightlist
\item[\circled{1}]
Diccionario con los valores actuales como claves y los valores por los
cuales serán reemplazados como valores.
\item[\circled{2}]
Hacemos el reemplazo.
\item[\circled{3}]
Convertimos a columna categórica.
\end{description}

\begin{verbatim}
<class 'pandas.core.frame.DataFrame'>
RangeIndex: 9 entries, 0 to 8
Data columns (total 3 columns):
 #   Column  Non-Null Count  Dtype   
---  ------  --------------  -----   
 0   A       9 non-null      int64   
 1   B       9 non-null      object  
 2   Acat    9 non-null      category
dtypes: category(1), int64(1), object(1)
memory usage: 485.0+ bytes
\end{verbatim}
\end{block}

\begin{block}{Categorizar una variable cuantitativa}
\phantomsection\label{categorizar-una-variable-cuantitativa}
Otra situación que se presenta con frecuencia es categorizar una
variable cuantitativa, por lo general de naturaleza continua. Por
ejemplo, tenemos los datos que se muestran en la tabla
\ref{tbl-datosedades} corresponden a edades de personas
\end{block}

\begin{block}{Datos}
\phantomsection\label{datos}
\begin{longtable}[]{@{}rrrrrrrrrr@{}}

\caption{\label{tbl-datosedades}Datos de edades de un grupo de personas}

\tabularnewline

\toprule\noalign{}
\endhead
\bottomrule\noalign{}
\endlastfoot
20 & 74 & 36 & 49 & 39 & 47 & 60 & 63 & 34 & 50 \\
46 & 44 & 28 & 28 & 39 & 65 & 52 & 65 & 55 & 73 \\
77 & 25 & 71 & 47 & 15 & 76 & 23 & 19 & 58 & 62 \\
57 & 25 & 59 & 46 & 39 & 42 & 25 & 21 & 19 & 55 \\
47 & 33 & 55 & 79 & 63 & 30 & 58 & 31 & 17 & 47 \\

\end{longtable}
\end{block}

\begin{block}{Categorizar la edad}
\phantomsection\label{categorizar-la-edad}
Usaremos los siguientes grupos etarios según la OMS: juventud {[}14 -
26{]} años, adultez joven (26 - 40{]} años, adultez intermedia (40 -
50{]} años, adultez tardía (50 - 60{]} años y adulto mayor 60 años o
más.

\phantomsection\label{annotated-cell-63}%
\begin{Shaded}
\begin{Highlighting}[]
\NormalTok{edad}\OperatorTok{=}\NormalTok{pd.Series([}\DecValTok{20}\NormalTok{, }\DecValTok{74}\NormalTok{, }\DecValTok{36}\NormalTok{, }\DecValTok{49}\NormalTok{, }\DecValTok{39}\NormalTok{, }\DecValTok{47}\NormalTok{, }\DecValTok{60}\NormalTok{, }\DecValTok{63}\NormalTok{, }\DecValTok{34}\NormalTok{, }\DecValTok{50}\NormalTok{,}
\DecValTok{46}\NormalTok{, }\DecValTok{44}\NormalTok{, }\DecValTok{28}\NormalTok{, }\DecValTok{28}\NormalTok{, }\DecValTok{39}\NormalTok{, }\DecValTok{65}\NormalTok{, }\DecValTok{52}\NormalTok{, }\DecValTok{65}\NormalTok{, }\DecValTok{55}\NormalTok{, }\DecValTok{73}\NormalTok{,}
\DecValTok{77}\NormalTok{, }\DecValTok{25}\NormalTok{, }\DecValTok{71}\NormalTok{, }\DecValTok{47}\NormalTok{, }\DecValTok{15}\NormalTok{, }\DecValTok{76}\NormalTok{, }\DecValTok{23}\NormalTok{, }\DecValTok{19}\NormalTok{, }\DecValTok{58}\NormalTok{, }\DecValTok{62}\NormalTok{,}
\DecValTok{57}\NormalTok{, }\DecValTok{25}\NormalTok{, }\DecValTok{59}\NormalTok{, }\DecValTok{46}\NormalTok{, }\DecValTok{39}\NormalTok{, }\DecValTok{42}\NormalTok{, }\DecValTok{25}\NormalTok{, }\DecValTok{21}\NormalTok{, }\DecValTok{19}\NormalTok{, }\DecValTok{55}\NormalTok{,}
\DecValTok{47}\NormalTok{, }\DecValTok{33}\NormalTok{, }\DecValTok{55}\NormalTok{, }\DecValTok{79}\NormalTok{, }\DecValTok{63}\NormalTok{, }\DecValTok{30}\NormalTok{, }\DecValTok{58}\NormalTok{, }\DecValTok{31}\NormalTok{, }\DecValTok{17}\NormalTok{, }\DecValTok{47}\NormalTok{] ) }\hspace*{\fill}\NormalTok{\circled{1}}
\NormalTok{edadf}\OperatorTok{=}\NormalTok{pd.cut(edad, bins}\OperatorTok{=}\NormalTok{[}\DecValTok{14}\NormalTok{,}\DecValTok{26}\NormalTok{,}\DecValTok{40}\NormalTok{,}\DecValTok{50}\NormalTok{,}\DecValTok{60}\NormalTok{,}\BuiltInTok{float}\NormalTok{(}\StringTok{\textquotesingle{}Inf\textquotesingle{}}\NormalTok{)],}
\NormalTok{labels}\OperatorTok{=}\NormalTok{ [}\StringTok{\textquotesingle{}Juventud\textquotesingle{}}\NormalTok{,}\StringTok{\textquotesingle{}Adultez joven\textquotesingle{}}\NormalTok{,}
\StringTok{\textquotesingle{}Adultez intermedia\textquotesingle{}}\NormalTok{,}\StringTok{\textquotesingle{}adultez tardía\textquotesingle{}}\NormalTok{,}\StringTok{\textquotesingle{}Adulto Mayor\textquotesingle{}}\NormalTok{]) }\hspace*{\fill}\NormalTok{\circled{2}}
\NormalTok{dfedad}\OperatorTok{=}\NormalTok{pd.DataFrame(\{}\StringTok{"Edad"}\NormalTok{:edad,}\StringTok{"Categ"}\NormalTok{:edadf\}) }\hspace*{\fill}\NormalTok{\circled{3}}
\NormalTok{dfedad.head() }\hspace*{\fill}\NormalTok{\circled{4}}
\end{Highlighting}
\end{Shaded}

\begin{description}
\tightlist
\item[\circled{1}]
Ingresamos los datos como una Serie de Pandas.
\item[\circled{2}]
Usamos la función \texttt{cut} \index{Pandas!Función!cut} para
categorizar.
\item[\circled{3}]
Colocamos en un marco de datos las edades y su categorización.
\item[\circled{4}]
Se imprimen las 5 primeras filas del marco de datos que contiene las
edades y su categorización.
\end{description}

\begin{verbatim}
   Edad               Categ
0    20            Juventud
1    74        Adulto Mayor
2    36       Adultez joven
3    49  Adultez intermedia
4    39       Adultez joven
\end{verbatim}
\end{block}

\begin{block}{}
\phantomsection\label{section-2}
Ahora podemos contar cuantas personas hay en cada grupo etáreo usando el
método \texttt{count()}. \index{Pandas Data Frame!Método!count}

\begin{Shaded}
\begin{Highlighting}[]
\NormalTok{edadf.value\_counts() }
\end{Highlighting}
\end{Shaded}

\begin{verbatim}
Adulto Mayor          11
Juventud              10
Adultez joven         10
Adultez intermedia    10
adultez tardía         9
Name: count, dtype: int64
\end{verbatim}

Estos métodos para contar, sumar, promediar y, en general, para calcular
estadísticas de resumen los estudiaremos en las siguientes clases.
\end{block}
\end{frame}




\end{document}
